% ============================================================================
% pygeoinf Technical Manual
% Theory, Algorithms, and API Reference
% ============================================================================
\documentclass[11pt,twoside,openright]{book}

% ---------------------------------------------------------------------------
% Typography and layout
% ---------------------------------------------------------------------------
\usepackage[final]{microtype}
\usepackage[T1]{fontenc}
\usepackage{lmodern}
\usepackage[margin=1in,headheight=14pt]{geometry}
\usepackage{fancyhdr}
\usepackage{titlesec}
\usepackage{parskip}          % paragraph spacing instead of indent

% ---------------------------------------------------------------------------
% Mathematics
% ---------------------------------------------------------------------------
\usepackage{amsmath,amssymb,amsthm}
\usepackage{mathtools}
% \usepackage{physics}  % install texlive-science if needed
\usepackage{bm}

% ---------------------------------------------------------------------------
% References, links, cross-references
% ---------------------------------------------------------------------------
\usepackage[hidelinks,bookmarksnumbered]{hyperref}
\usepackage[nameinlink,capitalise]{cleveref}
\usepackage[numbers,sort&compress]{natbib}

% ---------------------------------------------------------------------------
% Figures, diagrams, tables
% ---------------------------------------------------------------------------
\usepackage{graphicx}
\usepackage{tikz-cd}
\usepackage{tikz}
\usepackage{caption}
\usepackage{subcaption}
\usepackage{booktabs}
\usepackage{array}
\usepackage{longtable}

% ---------------------------------------------------------------------------
% Code and algorithm environments
% ---------------------------------------------------------------------------
\usepackage{listings}
\usepackage[ruled,vlined]{algorithm2e}

% ---------------------------------------------------------------------------
% Boxed callouts for code links
% ---------------------------------------------------------------------------
\usepackage[framemethod=tikz]{mdframed}

\newmdenv[
  linewidth=0.6pt,
  roundcorner=4pt,
  backgroundcolor=blue!3,
  linecolor=blue!30,
  skipabove=\medskipamount,
  skipbelow=\medskipamount,
  frametitle={\small\sffamily pygeoinf API},
  frametitlebackgroundcolor=blue!10,
  frametitlerule=true,
  frametitlerulewidth=0.4pt,
]{codebox}

\newmdenv[
  linewidth=0.6pt,
  roundcorner=4pt,
  backgroundcolor=yellow!5,
  linecolor=yellow!50!orange,
  skipabove=\medskipamount,
  skipbelow=\medskipamount,
  frametitle={\small\sffamily Why this matters},
  frametitlebackgroundcolor=yellow!15,
  frametitlerule=true,
  frametitlerulewidth=0.4pt,
]{whybox}

\newmdenv[
  linewidth=0.6pt,
  roundcorner=4pt,
  backgroundcolor=green!3,
  linecolor=green!40,
  skipabove=\medskipamount,
  skipbelow=\medskipamount,
  frametitle={\small\sffamily Worked Example},
  frametitlebackgroundcolor=green!10,
  frametitlerule=true,
  frametitlerulewidth=0.4pt,
]{examplebox}

\newmdenv[
  linewidth=0.6pt,
  roundcorner=4pt,
  backgroundcolor=red!3,
  linecolor=red!30,
  skipabove=\medskipamount,
  skipbelow=\medskipamount,
  frametitle={\small\sffamily Common Pitfall},
  frametitlebackgroundcolor=red!8,
  frametitlerule=true,
  frametitlerulewidth=0.4pt,
]{pitfallbox}

% ---------------------------------------------------------------------------
% Theorem environments
% ---------------------------------------------------------------------------
\theoremstyle{plain}
\newtheorem{theorem}{Theorem}[chapter]
\newtheorem{proposition}[theorem]{Proposition}
\newtheorem{lemma}[theorem]{Lemma}
\newtheorem{corollary}[theorem]{Corollary}

\theoremstyle{definition}
\newtheorem{definition}[theorem]{Definition}
\newtheorem{example}[theorem]{Example}

\theoremstyle{remark}
\newtheorem*{remark}{Remark}

% ---------------------------------------------------------------------------
% Custom commands: Spaces and maps
% ---------------------------------------------------------------------------
\newcommand{\modelspace}{\mathcal{M}}
\newcommand{\dataspace}{\mathcal{D}}
\newcommand{\propertyspace}{\mathcal{P}}
\newcommand{\Gd}{\mathbf{G}}              % standard discrete forward matrix
\newcommand{\GdGal}{\tilde{\mathbf{G}}}   % Galerkin discrete matrix
\newcommand{\Tau}{\mathcal{T}}             % property operator

% ---------------------------------------------------------------------------
% Custom commands: Measures
% ---------------------------------------------------------------------------
\newcommand{\prior}{\mu_{\modelspace}^0}
\newcommand{\post}{\mu_{\modelspace}^{\mathbf{\tilde{d}}}}
\newcommand{\discprior}{\mu_{\modelspace}^{0,N}}
\newcommand{\discpost}{\mu_{\modelspace}^{\mathbf{\tilde{d}}, N}}

% ---------------------------------------------------------------------------
% Custom commands: Discretization
% ---------------------------------------------------------------------------
\newcommand{\Vspace}{\mathbb{V}_N}
\newcommand{\mass}{\mathbf{M}}             % Gram/mass matrix

% ---------------------------------------------------------------------------
% Custom commands: Inner products and pairings
% ---------------------------------------------------------------------------
\newcommand{\ipM}[2]{(#1,#2)_{\modelspace}}
\newcommand{\ipV}[2]{(#1,#2)_{\Vspace}}
\newcommand{\ipD}[2]{(#1,#2)_{\dataspace}}
\newcommand{\ipH}[2]{(#1,#2)_{H}}
\newcommand{\ip}[3]{(#1,#2)_{#3}}         % generic inner product
\newcommand{\pair}[2]{\langle #1, #2\rangle}

% ---------------------------------------------------------------------------
% Custom commands: Riesz maps and adjoint
% ---------------------------------------------------------------------------
\newcommand{\RieszM}{\mathcal{R}_{\modelspace}}
\newcommand{\RieszD}{\mathcal{R}_{\dataspace}}
\newcommand{\Riesz}{\mathcal{R}}
\newcommand{\Gadj}{G^*}

% ---------------------------------------------------------------------------
% Custom commands: Math operators
% ---------------------------------------------------------------------------
\DeclareMathOperator*{\argmin}{arg\,min}
\DeclareMathOperator*{\argmax}{arg\,max}
\DeclareMathOperator{\spn}{span}
\DeclareMathOperator{\supp}{supp}
\DeclareMathOperator{\dom}{dom}
\DeclareMathOperator{\epi}{epi}
\DeclareMathOperator{\diag}{diag}
\DeclareMathOperator{\tr}{tr}
\DeclareMathOperator{\rank}{rank}
\DeclareMathOperator{\range}{range}
\DeclareMathOperator{\kernel}{ker}

% ---------------------------------------------------------------------------
% Listings style for Python
% ---------------------------------------------------------------------------
\lstdefinestyle{pygeoinf}{
  language=Python,
  basicstyle=\small\ttfamily,
  keywordstyle=\color{blue!70!black}\bfseries,
  commentstyle=\color{green!50!black}\itshape,
  stringstyle=\color{red!60!black},
  showstringspaces=false,
  breaklines=true,
  frame=none,
  xleftmargin=0.5em,
  aboveskip=0.5em,
  belowskip=0.5em,
}
\lstset{style=pygeoinf}

% ---------------------------------------------------------------------------
% Page headers and footers
% ---------------------------------------------------------------------------
\pagestyle{fancy}
\fancyhf{}
\fancyhead[LE]{\small\slshape\nouppercase{\leftmark}}
\fancyhead[RO]{\small\slshape\nouppercase{\rightmark}}
\fancyfoot[C]{\thepage}
\renewcommand{\headrulewidth}{0.4pt}

% ---------------------------------------------------------------------------
% Title information
% ---------------------------------------------------------------------------
\title{%
  \vspace{-2cm}
  {\Huge\bfseries pygeoinf}\\[0.8em]
  {\LARGE Technical Manual}\\[0.4em]
  {\Large Theory, Algorithms, and API Reference}\\[2em]
  \rule{\textwidth}{0.8pt}
}
\author{%
  Adrian\\[0.5em]
  {\normalsize Department of Earth Sciences}
}
\date{\today}

% ============================================================================
\begin{document}
% ============================================================================

\frontmatter

% ---------------------------------------------------------------------------
% Title page
% ---------------------------------------------------------------------------
\maketitle

% ---------------------------------------------------------------------------
% Preface
% ---------------------------------------------------------------------------
\chapter{Preface}

\section*{Who this manual is for}

This manual is written for geophysicists who have a forward operator and
a dataset and want to perform rigorous linear inversion or inference. The
word \emph{rigorous} carries weight here: pygeoinf does not cut corners.
It insists on correct Hilbert space structure, proper adjoints, and
treatment of uncertainty that is mathematically honest. If you are used
to inverting $( \mathbf{G}^\top \mathbf{G})^{-1}\mathbf{G}^\top$
and calling it done, this manual will explain why that is sometimes
wrong, and what to do instead.

No prior knowledge of functional analysis is assumed, but you will need
to engage with the mathematics. The package deliberately surfaces
conceptual machinery --- Riesz maps, dual pairings, Hilbert adjoints ---
that many inversion codes hide. The reward is a library whose results are
correct by construction for \emph{any} choice of model space, inner
product, and forward operator, not just the Euclidean special case.

\section*{How to read this manual}

Each chapter follows a consistent four-part structure:
\begin{enumerate}
  \item \emph{Why this matters} --- a motivating argument, usually with a
    geophysical example, that explains why the concept is worth the
    effort. Highlighted in yellow callout boxes.
  \item \emph{Definitions and theorems} --- the full, rigorous mathematical
    treatment.
  \item \emph{Implementation in pygeoinf} --- how each mathematical concept
    maps to the API. Highlighted in blue callout boxes labelled
    ``pygeoinf API''.
  \item \emph{Worked example} --- a complete mini-example, in both
    mathematics and code. Highlighted in green callout boxes.
\end{enumerate}

If you are primarily a practitioner, you can read the yellow and blue
boxes and skip the proofs on a first pass without losing the ability to
use the package.
If you are mathematically inclined, read everything; proofs that would
interrupt the flow are deferred to \cref{app:proofs}.

Red callout boxes labelled ``Common Pitfall'' flag mistakes that are
easy to make and important to avoid.

\section*{Acknowledgements}

This manual documents work carried out at the Department of Earth
Sciences. The theoretical framework draws heavily on
\citet{al2021linear}, \citet{stuart2010inverse}, and the trilogy
\citet{backus1970inference,backus1970uniqueness,backus1970inferenceIII}.

% ---------------------------------------------------------------------------
% Table of contents
% ---------------------------------------------------------------------------
\tableofcontents

% ---------------------------------------------------------------------------
% List of notation
% ---------------------------------------------------------------------------
\chapter{Notation and Conventions}
\label{notation}

\section*{General conventions}

Throughout this manual we use the following typographic conventions.

\begin{itemize}
  \item \textbf{Boldface roman} ($\mathbf{m}$, $\mathbf{G}$, $\boldsymbol{\lambda}$)
    denotes finite-dimensional vectors and matrices, i.e.\ elements of
    $\mathbb{R}^n$ or linear maps $\mathbb{R}^n \to \mathbb{R}^m$. We also use
    boldface for coefficient vectors/sequences and concrete operators acting on
    coefficient spaces (e.g.\ $\mathbf{m} \in \ell^2$ and $\mathbf{M}_\phi$).
  \item \textbf{Calligraphic} ($\mathcal{M}$, $\mathcal{D}$, $\mathcal{P}$)
    denotes abstract (possibly infinite-dimensional) Hilbert or Banach
    spaces.
  \item \textbf{Script/Roman capitals} ($G$, $T$, $\Pi$, $\pi$, $\mathcal{R}$)
    denote abstract linear operators between such spaces.
  \item Indices: $j,k,\ell$ typically index a basis of the model space;
    $i$ typically indexes a datum.
  \item We work exclusively over $\mathbb{R}$; all spaces are real.
\end{itemize}

\section*{Spaces}

\begin{longtable}{lll}
\toprule
Symbol & Meaning & Code \\
\midrule
$\modelspace$ & Model space (Hilbert) & \texttt{model\_space} \\
$\dataspace$ & Data space (usually $\mathbb{R}^{N_d}$) & \texttt{data\_space} \\
$\mathcal{P}$ & Property space (usually $\mathbb{R}^{N_p}$) & \texttt{property\_space} \\
$\modelspace'$ & Continuous dual of $\modelspace$ & implicit via Riesz \\
$\dataspace'$ & Continuous dual of $\dataspace$ & implicit via Riesz \\
$H$ & Generic Hilbert space & \texttt{HilbertSpace} \\
$\ell^2$ & Square-summable sequences & infinite-dim.~limit \\
$\mathbb{R}^N$ & Finite coefficient space & \texttt{EuclideanSpace(N)} \\
$\modelspace_N$ & $N$-dimensional subspace of $\modelspace$ & \texttt{dim = N} \\
$H_1 \oplus H_2$ & Hilbert direct sum & \texttt{HilbertSpaceDirectSum} \\
\bottomrule
\end{longtable}

\section*{Coefficient spaces and standard bases}

We use two closely related coefficient spaces:
\begin{itemize}
  \item the infinite-dimensional space $\ell^2$ of square-summable sequences, and
  \item its finite-dimensional truncation $\mathbb{R}^N$.
\end{itemize}
For $\ell^2$ we denote the standard basis vectors by $\{e_j\}_{j=1}^\infty$, where
$(e_j)_k = \delta_{jk}$. For $\mathbb{R}^N$ we denote the standard basis vectors by
$\{\mathbf{e}_j\}_{j=1}^N$.

On the dual coefficient spaces we use the corresponding \emph{coordinate functionals}
$\{e^j\}_{j=1}^\infty \subset (\ell^2)'$ and (optionally) $\{\mathbf{e}^j\}_{j=1}^N \subset (\mathbb{R}^N)'$, defined by
\[
  e^j(\mathbf{c}) = c_j, \qquad \mathbf{e}^j(\mathbf{a}) = a_j.
\]
They satisfy the canonical pairing relations
\[
  \langle e^i, e_j\rangle = \delta^i_j, \qquad \langle \mathbf{e}^i, \mathbf{e}_j\rangle = \delta^i_j.
\]

\begin{remark}
For the standard $\ell^2$ inner product $(\mathbf{c},\mathbf{d})_{\ell^2} := \sum_j c_j d_j$, the continuous dual satisfies $(\ell^2)' \cong \ell^2$ via the Riesz map; in particular, the coordinate functional $e^j$ corresponds to the vector $e_j$. We will nevertheless keep primes (e.g. $(\ell^2)'$) in the manual to make dual/primal roles explicit.
\end{remark}

\section*{Operators and maps}

\begin{longtable}{lll}
\toprule
Symbol & Meaning & Code \\
\midrule
$G : \modelspace \to \dataspace$ & Forward operator & \texttt{forward\_operator} \\
$G' : \dataspace' \to \modelspace'$ & Banach dual (transpose) & \texttt{.dual} \\
$G^* : \dataspace \to \modelspace$ & Hilbert adjoint & \texttt{.adjoint} \\
$\mathcal{T} : \modelspace \to \mathcal{P}$ & Property operator & \texttt{property\_operator} \\
$\mathcal{R}_X : X' \to X$ & Riesz isomorphism & \texttt{from\_dual} \\
$\mathcal{R}_X^{-1} : X \to X'$ & Inverse Riesz & \texttt{to\_dual} \\
$\Pi : \modelspace \to \ell^2$ & Analysis (expansion coefficients) & \texttt{to\_components} \\
$\pi : \ell^2 \to \modelspace$ & Synthesis (reconstruct from coeffs) & \texttt{from\_components} \\
$\Pi_N : \modelspace \to \mathbb{R}^N$ & Truncated analysis & \texttt{to\_components} (finite) \\
$\pi_N : \mathbb{R}^N \to \modelspace_N$ & Truncated synthesis & \texttt{from\_components} (finite) \\
$P_N : \modelspace \to \modelspace_N$ & Orthogonal projection onto $\modelspace_N$ & implicit \\
$\phi_j \in \modelspace$ & Basis vector & \texttt{basis\_vector(j)} \\
$\psi_j \in \modelspace$ & Biorthogonal dual basis vector & derived from Gram \\
\bottomrule
\end{longtable}

\section*{Inner products and pairings}

\begin{longtable}{lll}
\toprule
Symbol & Meaning & Code \\
\midrule
$(m_1, m_2)_{\modelspace}$ & Inner product on $\modelspace$ & \texttt{inner\_product(x1,x2)} \\
$(d_1, d_2)_{\dataspace}$ & Inner product on $\dataspace$ & \texttt{inner\_product(d1,d2)} \\
$\langle \ell, m \rangle$ & Dual pairing: $\ell \in \modelspace',\, m \in \modelspace$ & \texttt{l(m)} \\
$\|m\|_{\modelspace}$ & Hilbert norm $\sqrt{(m,m)_{\modelspace}}$ & \texttt{norm(m)} \\
\bottomrule
\end{longtable}

We use \emph{parentheses} $(\cdot,\cdot)$ for inner products and
\emph{angle brackets} $\langle\cdot,\cdot\rangle$ for dual pairings.
On a Hilbert space the two are related by
$\langle \ell, m \rangle = (m,\, \mathcal{R}_X \ell)_X$.

\section*{Matrices (finite-dimensional objects)}

\begin{longtable}{lll}
\toprule
Symbol & Meaning & Notes \\
\midrule
$\mathbf{M}_{\phi,N}$ & Gram/mass matrix $[\,(\phi_k,\phi_j)_{\modelspace}\,]_{j,k}$ & SPD, $N\times N$ \\
$\mathbf{G}$ & Standard discrete matrix of $G$: $[\,(G\phi_j,\beta_i)_{\dataspace}\,]$ & $N_d \times N$ \\
$\tilde{\mathbf{G}}$ & Galerkin matrix (with inner products) & \texttt{.matrix(galerkin=True)} \\
$\mathbf{m}, \mathbf{d}$ & Coefficient vectors in $\mathbb{R}^N$, $\mathbb{R}^{N_d}$ & \texttt{to\_components(m)} \\
\bottomrule
\end{longtable}

\section*{Measures and probability}

\begin{longtable}{lll}
\toprule
Symbol & Meaning & Code \\
\midrule
$\mu$ & A probability measure on $\modelspace$ & \\
$\mathcal{N}(\mu_0, C)$ & Gaussian measure, mean $\mu_0$, covariance $C$ & \texttt{GaussianMeasure} \\
$\mu^0$ & Prior measure & \texttt{prior} \\
$\mu^{\tilde{d}}$ & Posterior measure given data $\tilde{d}$ & \texttt{.invert(data)} \\
$C : \modelspace \to \modelspace$ & Covariance operator & \texttt{.covariance\_operator} \\
$\mathbf{C}_{\dataspace} : \dataspace \to \dataspace$ & Data noise covariance & \texttt{data\_error\_measure} \\
$\varepsilon \in \dataspace$ & Data noise, $\varepsilon \sim \mathcal{N}(0,\mathbf{C}_{\dataspace})$ & \\
\bottomrule
\end{longtable}

\section*{Convex analysis}

\begin{longtable}{lll}
\toprule
Symbol & Meaning & Code \\
\midrule
$C \subset \modelspace$ & A convex subset & \texttt{Subset} subclass \\
$h_C(q) = \sigma_C(q)$ & Support function of $C$ & \texttt{SupportFunction} \\
$\partial h_C(q)$ & Subdifferential of $h_C$ at $q$ & \texttt{.subgradient(q)} \\
$\mathcal{U} \subset \mathcal{P}$ & Admissible property set & computed by DLI \\
\bottomrule
\end{longtable}

% ============================================================================
\mainmatter
% ============================================================================

% %%%%%%%%%%%%%%%%%%%%%%%%%%%%%%%%%%%%%%%%%%%%%%%%%%%%%%%%%%%%%%%%%%%%%%%%%%%%
\part{Foundations}
\label{part:foundations}
% %%%%%%%%%%%%%%%%%%%%%%%%%%%%%%%%%%%%%%%%%%%%%%%%%%%%%%%%%%%%%%%%%%%%%%%%%%%%

% ============================================================================
\chapter{Introduction and Motivation}
\label{ch:introduction}
% ============================================================================

\section{What is pygeoinf?}
\label{sec:what-is-pygeoinf}

pygeoinf is a Python library for rigorous linear geophysical inversion and
Bayesian inference. It is built around a single design philosophy:
\emph{the mathematical structure of the problem should be encoded explicitly
in the code}, not reduced to bare matrices and vectors.

Most inversion codes operate on the discrete level from the start: the
user assembles a matrix $\mathbf{G}\in\mathbb{R}^{M\times N}$, a data
vector $\mathbf{d}\in\mathbb{R}^M$, and calls some solver. This approach
works, but it forces the user to manually track which inner products
are implicit, which transposes are adjoints, and how uncertainty
propagates through the discretization. When anything non-standard
happens --- a non-Euclidean norm, a non-uniform grid, a Sobolev
prior --- the bookkeeping becomes error-prone.

Instead, pygeoinf exposes the following conceptual objects as first-class
Python types:

\begin{itemize}
  \item \textbf{Hilbert spaces} (\texttt{HilbertSpace}) --- carrying their
    inner product, Riesz maps, and basis structure.
  \item \textbf{Linear operators} (\texttt{LinearOperator}) --- aware of
    their domain and codomain, automatically computing the correct Hilbert
    adjoint from the spaces' Riesz maps.
  \item \textbf{Linear forms} (\texttt{LinearForm}) --- elements of the dual
    space, representing individual measurements.
  \item \textbf{Gaussian measures} (\texttt{GaussianMeasure}) --- priors and
    posteriors as probability measures on Hilbert spaces.
  \item \textbf{Convex sets and support functions} --- geometry of the
    admissible model set.
\end{itemize}

Once these objects are defined, a full suite of inversion algorithms ---
Bayesian, least-squares, minimum-norm, Backus-Gilbert, and
deterministic linear inference --- work automatically and correctly
for any choice of spaces and operators.

\section{The geophysical inversion problem}
\label{sec:geophysical-inversion}

The starting point is a physical system described by a \emph{model}
$m$ --- perhaps the seismic wave-speed field $v_P(x)$, the electrical
conductivity $\sigma(x)$, or the subsurface density $\rho(x)$. Some
aspects of this system are measurable: we collect \emph{data}
$\tilde{d} = (\tilde{d}_1,\ldots,\tilde{d}_M)$ via instruments
(travel-time picks, gravity readings, magnetic measurements, \ldots).
A \emph{forward model} $G$ maps a model to its predicted data:
\begin{equation}
  \label{eq:forward-model}
  d = Gm + \varepsilon,
\end{equation}
where $\varepsilon$ is measurement noise. The \emph{inverse problem} is
to determine what equation~\eqref{eq:forward-model} implies about
$m$ given the observed data $\tilde{d}$.

All realistic geophysical inverse problems share three features:

\begin{enumerate}
  \item \textbf{Underdetermination.} The data $\tilde{d}$ is a
    finite-dimensional vector, whereas $m$ is a function (infinite-dimensional
    in principle). Infinitely many models are consistent with the data
    exactly; the problem has a non-trivial null space.
  \item \textbf{Ill-conditioning.} Even the components of $m$ that \emph{are}
    constrained by the data may be constrained very weakly, amplifying
    noise in the solution.
  \item \textbf{Incomplete prior knowledge.} We have partial information
    about $m$ --- physical bounds, smoothness expectations, perhaps a
    statistical prior --- that should inform the inversion.
\end{enumerate}

Any honest inversion framework must confront all three. pygeoinf does so
by working at the level of \emph{abstract spaces and operators}, which
makes the structure of the problem explicit and prevents silent,
basis-dependent errors.

\section{Why abstract spaces matter}
\label{sec:why-spaces-matter}

\begin{whybox}
If you are used to working with matrices and vectors in NumPy, you may
wonder why pygeoinf insists on Hilbert spaces, dual pairings, and Riesz
maps. This section shows you exactly what goes wrong when you ignore
these structures --- and how pygeoinf gets it right automatically.
\end{whybox}

To make the argument concrete, consider a simple 1D problem. You want to
recover a function $m : [0,1] \to \mathbb{R}$ from $M$ integral
measurements
\begin{equation}
  \label{eq:motivating-forward}
  d_i = \int_0^1 g_i(x)\, m(x)\, \mathrm{d}x, \qquad i = 1,\ldots,M.
\end{equation}
You discretize $[0,1]$ into $N$ cells with centres $x_j^*$ and widths
$\Delta x_j > 0$ (possibly non-uniform). Approximate $m$ as piecewise
constant: $m(x) \approx m_j$ on cell $j$, so $\mathbf{m} =
(m_1,\ldots,m_N)^\top \in \mathbb{R}^N$. The discrete forward matrix
becomes
\begin{equation}
  \label{eq:motivating-matrix}
  G_{ij} = g_i(x_j^*) \Delta x_j,
  \qquad \mathbf{d} = \mathbf{G}\mathbf{m} \in \mathbb{R}^M.
\end{equation}
You want the minimum-norm solution: among all $\mathbf{m}$ with
$\mathbf{G}\mathbf{m} = \mathbf{d}$, find the one of smallest
norm.

\paragraph{The naive approach: use the Euclidean norm on $\mathbb{R}^N$.}
Minimizing $\|\mathbf{m}\|_2^2 = \sum_j m_j^2$ gives
\begin{equation}
  \label{eq:naive-min-norm}
  \mathbf{m}^\dagger_{\text{naive}}
  = \mathbf{G}^\top (\mathbf{G}\mathbf{G}^\top)^{-1} \mathbf{d}.
\end{equation}

\paragraph{The correct approach: use the $L^2([0,1])$ norm.}
The natural inner product approximating $\|m\|_{L^2}^2 = \int_0^1 m(x)^2
\,\mathrm{d}x$ is
\begin{equation}
  \label{eq:L2-inner-product}
  (\mathbf{m},\mathbf{n})_{L^2} \approx \sum_j m_j n_j \Delta x_j
  = \mathbf{m}^\top \mathbf{M} \mathbf{n},
  \qquad \mathbf{M} = \mathrm{diag}(\Delta x_1,\ldots,\Delta x_N).
\end{equation}
Minimizing this gives the qualitatively different solution
\begin{equation}
  \label{eq:correct-min-norm}
  \mathbf{m}^\dagger
  = \mathbf{M}^{-1}\mathbf{G}^\top
    (\mathbf{G}\mathbf{M}^{-1}\mathbf{G}^\top)^{-1} \mathbf{d}.
\end{equation}

\paragraph{Why they differ.}
If the grid is uniform ($\Delta x_j = 1/N$ for all $j$), then
$\mathbf{M} = \tfrac{1}{N}\mathbf{I}$ and the two solutions agree up to
a constant. But on a non-uniform grid, the naive solution
\eqref{eq:naive-min-norm} \emph{implicitly weights finely-sampled
regions more heavily}: minimizing $\sum_j m_j^2$ without the $\Delta x_j$
weights is not minimizing a physically meaningful norm of $m$.

The matrix $\mathbf{M}$ is the \emph{Gram matrix} (or mass matrix) of the
basis $\{\phi_j\}$ where $\phi_j = \mathbf{1}_{\text{cell }j} / \Delta x_j^{1/2}$.
The formula $G^* = \mathbf{M}^{-1}\mathbf{G}^\top$ is the finite-dimensional
instance of the abstract relation
\begin{equation}
  \label{eq:hilbert-adjoint-abstract}
  G^* = \mathcal{R}_{\modelspace} \circ G' \circ \mathcal{R}_{\dataspace}^{-1},
\end{equation}
which pygeoinf enforces automatically once you provide the spaces with
their correct inner products.

\begin{pitfallbox}
The matrix transpose $\mathbf{G}^\top$ is the adjoint only in Euclidean
space with the standard dot product. In any other Hilbert space it is
not. Using $\mathbf{G}^\top$ instead of
$\mathbf{M}_{\mathcal{M}}^{-1}\mathbf{G}^\top\mathbf{M}_{\mathcal{D}}$
produces a solution that is biased toward densely-gridded regions (or,
in a Bayesian setting, posterior covariances that are not
basis-independent).
\end{pitfallbox}

This example is the simplest case. As problems become more complex ---
Sobolev regularity norms, non-trivial data weighting, operators between
function spaces --- the gap between ``naive'' and ``correct'' widens.
pygeoinf closes that gap structurally, not by asking the user to remember
the right formula each time.

\section{Package architecture overview}
\label{sec:architecture}

pygeoinf is organized in four ascending layers.

\medskip
\begin{center}
\begin{tikzpicture}[
  layer/.style={draw, rounded corners=4pt, fill=#1, minimum width=8cm,
                minimum height=0.75cm, font=\small\sffamily},
  arr/.style={->, thick, gray!60}
]
  \node[layer=blue!10]   (L1) at (0,0)    {Foundation: spaces, operators, linear forms};
  \node[layer=green!10]  (L2) at (0,1.0)  {Geometry: subsets, subspaces, support functions, measures};
  \node[layer=yellow!15] (L3) at (0,2.0)  {Algorithms: forward problems, inversions, optimisation};
  \node[layer=red!8]     (L4) at (0,3.0)  {Output: visualisation, diagnostics};
  \draw[arr] (L1) -- (L2);
  \draw[arr] (L2) -- (L3);
  \draw[arr] (L3) -- (L4);
\end{tikzpicture}
\end{center}
\medskip

\paragraph{Foundation layer}
(\cref{part:foundations}) The abstract base classes
\texttt{HilbertSpace}, \texttt{LinearOperator}, and \texttt{LinearForm}
define the mathematical objects.  Everything else is built on top of
them.  No inversion algorithm is aware of the concrete space type; it
only calls the abstract interface.

\paragraph{Geometry layer}
(\cref{part:geometry}) Convex subsets (\texttt{Ball},
\texttt{PolyhedralSet}, \ldots), subspaces (\texttt{AffineSubspace}),
support functions (\texttt{SupportFunction}), and Gaussian measures
(\texttt{GaussianMeasure}) all live here.  They each accept any
\texttt{HilbertSpace} and work correctly regardless of its concrete type.

\paragraph{Algorithm layer}
(\cref{part:algorithms}) \texttt{ForwardProblem},
\texttt{LinearBayesianInversion}, \texttt{LinearMinimumNormInversion},
\texttt{BackusGilbertInversion}, and the convex optimisation algorithms
belong here.  They compose the foundation and geometry objects into
complete solution procedures.

\paragraph{Output layer}
(\cref{part:computation}) \texttt{SubspaceSlicePlotter} and related
tools for visualising results and checking algorithmic correctness.

\section{How to read this manual}
\label{sec:how-to-read}

The manual is divided into four parts mirroring the package architecture:

\begin{description}
  \item[\cref{part:foundations} (Foundations)]
    Chapters~\ref{ch:hilbert-spaces}--\ref{ch:nonlinear}.
    The mathematical vocabulary: spaces, operators, forms.
    \emph{All users should read at least the motivation sections of these
    chapters.}
  \item[\cref{part:geometry} (Geometry)]
    Chapters~\ref{ch:subsets}--\ref{ch:gaussian-measures}.
    Convex geometry, subspaces, and probabilistic models.
  \item[\cref{part:algorithms} (Problems and Algorithms)]
    Chapters~\ref{ch:forward-problems}--\ref{ch:convex-optimisation}.
    How to set up and solve inverse problems.
  \item[\cref{part:computation} (Computational Methods)]
    Chapters~\ref{ch:computational}--\ref{ch:visualisation}.
    Numerical backends and visualisation.
\end{description}

Each chapter is self-contained. Forward references are kept to a minimum.
The appendices collect a convex analysis primer
(\cref{app:convex-primer}), functional analysis prerequisites
(\cref{app:functional-analysis}), deferred proofs
(\cref{app:proofs}), and a complete API quick-reference
(\cref{app:api-reference}).

\paragraph{Callout boxes.}
\begin{description}
  \item[\textcolor{yellow!60!gray}{Yellow} --- Why this matters]
    Motivational content and geophysical perspective.
  \item[\textcolor{blue!60}{Blue} --- pygeoinf API]
    The mapping from mathematical concept to Python code.
  \item[\textcolor{green!60!gray}{Green} --- Worked Example]
    Complete mini-examples with both equations and code.
  \item[\textcolor{red!60}{Red} --- Common Pitfall]
    Frequent mistakes and how to avoid them.
\end{description}

\section{Notation and conventions summary}
\label{sec:notation-summary}

A complete notation table is given in the \hyperref[notation]{Notation
and Conventions} chapter of the front matter. The most important rules:

\begin{itemize}
  \item Parentheses $(m_1, m_2)_{\modelspace}$ always denote inner products;
    angle brackets $\langle \ell, m \rangle$ always denote dual pairings.
  \item Boldface ($\mathbf{m}$, $\mathbf{G}$) is reserved for
    finite-dimensional vectors and matrices.
  \item Calligraphic letters ($\modelspace$, $\dataspace$) denote abstract
    spaces; Roman capitals ($G$, $T$, $\Pi$) denote abstract operators.
  \item $G^*$ always means the \emph{Hilbert} adjoint
    ($G^*: \dataspace \to \modelspace$); $G'$ means the \emph{Banach} dual
    ($G': \dataspace' \to \modelspace'$). They are related by
    \eqref{eq:hilbert-adjoint-abstract} but are not the same map.
\end{itemize}


% ============================================================================
\chapter{Hilbert Spaces}
\label{ch:hilbert-spaces}
% ============================================================================

\section{Why Hilbert spaces?}
\label{sec:why-hilbert}

\begin{whybox}
The inner product on your model space is not a minor technical detail ---
it determines what ``similar models'' means, and therefore what your
inversion algorithm considers a good solution. Two inversions that differ
only in their choice of inner product will in general produce qualitatively
different results. pygeoinf makes the inner product an explicit,
first-class object so that this choice is always visible and always
consistent.
\end{whybox}

Why not work with plain vectors and matrices? A vector space gives you
addition and scalar multiplication, but it has no notion of length or
angle. You cannot define a ``smallest'' element, orthogonality, or
projection without more structure. The extra structure needed is an
\emph{inner product}.

For geophysical inversion, the inner product on the model space $\modelspace$
is the mathematical encoding of your regularity expectations about the
model. A Euclidean inner product treats all degrees of freedom equally,
regardless of grid spacing or physical units. An $L^2([0,1])$ inner
product weights by area element $dx$. A Sobolev $H^1$ inner product
additionally penalises gradients. Each choice leads to a different
minimum-norm solution; only the ``right'' one is physically meaningful.

In short: \emph{the inner product is part of the problem specification},
not an implementation detail.

\section{Definition and axioms}
\label{sec:hilbert-definition}

\begin{definition}[Inner product space]
\label{def:inner-product-space}
A \emph{real inner product space} is a real vector space $H$ together with
a map $(\cdot,\cdot)_H : H \times H \to \mathbb{R}$ satisfying, for all
$u, v, w \in H$ and $\alpha \in \mathbb{R}$:
\begin{enumerate}
  \item \textbf{Symmetry:} $(u, v)_H = (v, u)_H$.
  \item \textbf{Linearity in the first argument:}
    $(\alpha u + v, w)_H = \alpha(u, w)_H + (v, w)_H$.
  \item \textbf{Positive definiteness:} $(u, u)_H \geq 0$, with equality only
    if $u = 0$.
\end{enumerate}
\end{definition}

The inner product induces a norm $\|u\|_H := \sqrt{(u, u)_H}$ and hence a
metric $d(u, v) = \|u - v\|_H$.

\begin{definition}[Hilbert space]
\label{def:hilbert-space}
A real \emph{Hilbert space} is a real inner product space $(H, (\cdot,\cdot)_H)$
that is \emph{complete} with respect to the induced norm: every Cauchy
sequence converges to an element of $H$.
\end{definition}

\begin{remark}
pygeoinf works exclusively over $\mathbb{R}$. Complex Hilbert spaces are not
supported. All spaces are assumed to be \emph{separable}, meaning they
admit a countable dense subset; this ensures the existence of countable
orthonormal bases.
\end{remark}

\paragraph{Orthogonality.}
Two elements $u, v \in H$ are \emph{orthogonal}, written $u \perp v$, if
$(u, v)_H = 0$. A set $S \subset H$ is orthogonal if any two distinct
elements of $S$ are orthogonal. If additionally $\|s\|_H = 1$ for all
$s \in S$, the set is \emph{orthonormal}.

\paragraph{The parallelogram law.}
A norm comes from an inner product if and only if it satisfies the
\emph{parallelogram law}:
\begin{equation}
  \label{eq:parallelogram}
  \|u + v\|^2_H + \|u - v\|^2_H = 2\|u\|^2_H + 2\|v\|^2_H
  \qquad \forall\, u, v \in H.
\end{equation}
The $\ell^p$ and $L^p$ spaces for $p \neq 2$ are Banach but not Hilbert;
they do not satisfy \eqref{eq:parallelogram}.

\section{Inner products, norms, and distances}
\label{sec:inner-products-norms}

\begin{proposition}[Cauchy--Schwarz inequality]
\label{prop:cauchy-schwarz}
For any $u, v \in H$,
\begin{equation}
  \label{eq:cauchy-schwarz}
  |(u, v)_H| \leq \|u\|_H \, \|v\|_H,
\end{equation}
with equality if and only if $u$ and $v$ are linearly dependent.
\end{proposition}

The Cauchy--Schwarz inequality guarantees that the cosine of the
``angle'' between two vectors, $\cos\theta_{u,v} := (u,v)_H / (\|u\|_H
\|v\|_H)$, is well-defined and in $[-1,1]$, even in infinite dimensions.

\paragraph{Polarisation identity.}
The inner product can be recovered from the norm via
\begin{equation}
  \label{eq:polarisation}
  (u, v)_H = \tfrac{1}{4}\left(\|u + v\|^2_H - \|u - v\|^2_H\right).
\end{equation}
This means all metric information about a Hilbert space is encoded in
its inner product.

\begin{codebox}
\texttt{HilbertSpace} exposes:
\begin{itemize}
  \item \texttt{inner\_product(x1, x2)}: computes $(x_1, x_2)_H$;
  \item \texttt{norm(x)}: computes $\|x\|_H$;
  \item \texttt{squared\_norm(x)}: computes $\|x\|^2_H$ (often cheaper);
  \item \texttt{distance(x1, x2)}: computes $\|x_1 - x_2\|_H$.
\end{itemize}
Internally, \texttt{inner\_product} is computed via the dual pairing
(see \cref{sec:dual-space}) as $(x_1, x_2)_H = \langle \mathcal{R}^{-1} x_1,\, x_2\rangle$,
using \texttt{to\_dual} and the duality product.
\end{codebox}

\section{The dual space and dual pairing}
\label{sec:dual-space}

\begin{whybox}
In geophysical terms, a dual vector is exactly what an instrument
measures: a number that depends linearly on the model. Understanding the
distinction between vectors (models) and dual vectors (measurements) is
fundamental to understanding why inversion requires Riesz maps.
\end{whybox}

\begin{definition}[Continuous dual space]
\label{def:dual-space}
Let $H$ be a Hilbert space. A \emph{continuous linear functional} on $H$
is a map $\ell : H \to \mathbb{R}$ that is linear ($\ell(\alpha u + v) =
\alpha \ell(u) + \ell(v)$) and bounded ($|\ell(u)| \leq C\|u\|_H$ for
some $C < \infty$). The collection of all continuous linear functionals
on $H$ forms the \emph{(continuous) dual space} $H'$, which is itself a
Banach space under the operator norm $\|\ell\|_{H'} := \sup_{\|u\|=1}
|\ell(u)|$.
\end{definition}

\paragraph{Dual pairing.}
We denote the evaluation of $\ell \in H'$ at $u \in H$ by the \emph{dual
pairing}:
\begin{equation}
  \label{eq:dual-pairing}
  \langle \ell, u \rangle := \ell(u).
\end{equation}
We use angle brackets $\langle\cdot,\cdot\rangle$ for dual pairings and
parentheses $(\cdot,\cdot)_H$ for inner products throughout this manual;
the two are related by the Riesz theorem (\cref{sec:riesz-theorem}).

\paragraph{Two distinct spaces.}
Elements of $H$ are ``primal'' vectors (models); elements of $H'$ are
``dual'' vectors (measurements, functionals, gradients of scalar
functions). On a finite-dimensional Euclidean space, $H \cong H'$ via
the standard dot product, which is why the distinction is often ignored
in numerical practice. On a general Hilbert space --- or more acutely, on
a Banach space where no inner product exists --- the two are different
objects and must be handled separately.

\paragraph{Concrete example.}
Let $H = L^2([0,1])$ and let $g \in L^2([0,1])$. Then the map $\ell:
m \mapsto \int_0^1 g(x) m(x)\,\mathrm{d}x$ is a continuous linear
functional on $H$, i.e.\ $\ell \in H'$. This is precisely the structure
of a geophysical measurement: every seismogram, gravity reading, or
magnetic datum is a dual vector.

\section{The Riesz representation theorem}
\label{sec:riesz-theorem}

\begin{theorem}[Riesz representation]
\label{thm:riesz}
Let $H$ be a real Hilbert space. For every continuous linear functional
$\ell \in H'$ there exists a unique $u_\ell \in H$ such that
\begin{equation}
  \label{eq:riesz-rep}
  \ell(v) = (u_\ell,\, v)_H \qquad \forall\, v \in H.
\end{equation}
Moreover, $\|u_\ell\|_H = \|\ell\|_{H'}$. The map $\Riesz: H' \to H$
defined by $\Riesz(\ell) := u_\ell$ is an isometric isomorphism.
\end{theorem}

\paragraph{Proof sketch.}
If $\ell = 0$ take $u_\ell = 0$. Otherwise, $\ker\ell$ is a closed
proper subspace of $H$ of codimension one. By the projection theorem
there exists a unit vector $e \perp \ker\ell$. Set $u_\ell :=
\ell(e)\,e$. For any $v \in H$, the element $v - \langle\ell,v\rangle e /
\ell(e)$ lies in $\ker\ell$, hence is orthogonal to $e$. A brief
calculation shows $(u_\ell, v)_H = \ell(v)$. Uniqueness follows from
positive definiteness.

\paragraph{Geometric interpretation.}
The Riesz representative $u_\ell = \Riesz(\ell)$ is the element of $H$
that ``points in the direction'' the functional $\ell$ is measuring. In
optimisation this is the \emph{gradient}: if $f(u) = \ell(u)$ then
$\nabla f = u_\ell$.

\paragraph{Riesz map and its inverse.}
We call $\Riesz : H' \to H$ the \emph{Riesz isomorphism} or Riesz map.
Its inverse $\Riesz^{-1} : H \to H'$ is defined by $(\Riesz^{-1} u)(v) =
(u, v)_H$. Together they establish the identification
\begin{equation}
  \label{eq:riesz-iden}
  \langle \ell, v \rangle = (v,\; \Riesz\,\ell)_H
  \qquad \forall\; \ell \in H',\; v \in H.
\end{equation}

\begin{codebox}
The Riesz isomorphism and its inverse are implemented as:
\[
  \texttt{from\_dual} \equiv \Riesz : H' \to H, \qquad
  \texttt{to\_dual}   \equiv \Riesz^{-1} : H \to H'.
\]
If you need them as \texttt{LinearOperator} instances, use the properties
\texttt{HilbertSpace.riesz} ($\Riesz$) and
\texttt{HilbertSpace.inverse\_riesz} ($\Riesz^{-1}$). The dual pairing
evaluates as $\langle \ell, v\rangle = \ell(v)$, implemented by calling
the \texttt{LinearForm} object directly: \texttt{l(v)}.
\end{codebox}

\section{Riesz-first design in pygeoinf}
\label{sec:riesz-first-design}

The key design choice in \texttt{pygeoinf} is that the \emph{primitive} piece of
Hilbert-space structure is not an inner-product formula, but the pair of maps
implementing the Riesz identification:
\(\Riesz : H' \to H\) and \(\Riesz^{-1} : H \to H'\), exposed in the API as
\texttt{from\_dual} and \texttt{to\_dual}.
Once these are defined, the inner product is \emph{derived} via the duality
pairing.

\paragraph{Inner products from duality.}
Given $u,v \in H$, we define
\begin{equation}
  \label{eq:ip-from-to-dual}
  (u,v)_H \,=\, \pair{\Riesz^{-1}u}{v}
  \,=\, \pair{\texttt{to\_dual}(u)}{v}.
\end{equation}
In code this is the default implementation
\begin{equation*}
  \texttt{inner\_product(u,v)} \,=\, \texttt{duality\_product(to\_dual(u), v)}.
\end{equation*}
Consequently, norms and distances are also derived:
\(\|u\|_H = \sqrt{(u,u)_H}\) and \(\|u-v\|_H\) uses the same geometry.

\paragraph{Adjoints depend on the geometry.}
For a bounded linear operator $A : H_1 \to H_2$, the Hilbert adjoint
$A^* : H_2 \to H_1$ is characterised by
\begin{equation}
  \label{eq:adjoint-characterisation}
  (A u,\, v)_{H_2} = (u,\, A^* v)_{H_1}
  \qquad \forall\; u \in H_1,\; v \in H_2.
\end{equation}
Because the inner products on $H_1$ and $H_2$ are defined through their Riesz
maps, the adjoint is inherently geometry-dependent. In \textit{pygeoinf}, when an
operator's adjoint is \emph{derived by the default machinery} (from the dual map
and the Riesz identifications), the property \texttt{LinearOperator.adjoint}
depends on \texttt{to\_dual}/\texttt{from\_dual} of the operator's domain and
codomain. Changing those maps therefore changes the induced geometry (norms,
gradients, and these default adjoints) in a consistent way across the API.
If an operator provides an explicit \texttt{adjoint\_mapping}, that
implementation takes precedence and can override the Riesz-derived default.

\begin{codebox}
Notation \(\leftrightarrow\) \texttt{pygeoinf} API (Riesz-first viewpoint):
\[
  \Riesz \equiv \texttt{HilbertSpace.from\_dual}, \qquad
  \Riesz^{-1} \equiv \texttt{HilbertSpace.to\_dual}.
\]
\[
  (u,v)_H \equiv \texttt{HilbertSpace.inner\_product(u,v)}
  \;=\; \pair{\texttt{to\_dual}(u)}{v}.
\]
\[
  \texttt{HilbertSpace.riesz} \equiv \Riesz, \qquad
  \texttt{HilbertSpace.inverse\_riesz} \equiv \Riesz^{-1}.
\]
Adjoints are geometry-aware by default: when \texttt{LinearOperator.adjoint} is
derived (rather than explicitly provided), it uses the domain/codomain Riesz
maps to satisfy \eqref{eq:adjoint-characterisation}. An explicitly-provided
\texttt{adjoint\_mapping} overrides this derivation.
\end{codebox}


\section{Bases and Representations: The Primal Side}

In numerical work we may not always want to manipulate an abstract element
$m \in H$ directly.
Instead, we might want to work with a sequence of numbers --- its \emph{coefficients}.
To make this precise, we must choose a basis.

Let $H$ be a Hilbert space with inner product $(\cdot,\cdot)_H$,
and let $\{\phi_j\}_{j=1}^\infty \subset H$ be a Riesz basis.

By definition, there exist constants $0 < A \le B < \infty$ such that
for every finitely supported coefficient sequence $\mathbf{c} \in \ell^2$,
\begin{equation}
A \|\mathbf{c}\|_{\ell^2}^2
\le
\left\|\sum_{j} [\mathbf{c}]_j \phi_j\right\|_H^2
\le
B \|\mathbf{c}\|_{\ell^2}^2 .
\label{eq:riesz-bounds-ped}
\end{equation}

This inequality has an important interpretation:

\begin{itemize}
\item The vector $\sum_j [\mathbf{c}]_j \phi_j$ depends continuously on the coefficients.
\item No information is lost or artificially amplified.
\item Coefficient vectors and physical vectors measure size in comparable ways.
\end{itemize}

In other words, the basis is \emph{stable}.

\subsection{Synthesis: From Coefficients to Vectors}

Given a sequence of coefficients $\mathbf{c} = ([\mathbf{c}]_j)$, we construct a vector in $H$ by
\begin{equation}
\pi(\mathbf{c}) := \sum_{j=1}^\infty [\mathbf{c}]_j \phi_j .
\end{equation}

This defines the \emph{synthesis operator}
\[
\pi : \ell^2 \to H.
\]

The Riesz bounds \eqref{eq:riesz-bounds-ped} imply:

\begin{itemize}
\item $\pi$ is linear,
\item $\pi$ is bounded,
\item $\pi$ is bijective,
\item and its inverse is bounded.
\end{itemize}

Thus $\pi$ is a topological isomorphism between $\ell^2$ and $H$.
This justifies thinking of $\ell^2$ as the \emph{coefficient space}
associated with the basis $\{\phi_j\}$.

\subsection{Analysis: From Vectors to Coefficients}

Since $\pi$ is bijective, each vector $m \in H$ has a unique
coefficient sequence $\mathbf{c}$ satisfying $\pi(\mathbf{c}) = m$.

This defines the \emph{analysis operator}
\[
\Pi := \pi^{-1} : H \to \ell^2.
\]

To compute coefficients explicitly, we introduce the
\emph{biorthogonal dual basis} $\{\psi_j\} \subset H$,
defined by
\begin{equation}
(\phi_j,\psi_k)_H = \delta_{jk}.
\end{equation}

Then the coefficients of $m$ are given by
\begin{equation}
\Pi(m) = \bigl[(m,\psi_j)_H\bigr]_j.
\end{equation}

Thus:

\[
\pi \circ \Pi = \mathrm{Id}_H,
\qquad
\Pi \circ \pi = \mathrm{Id}_{\ell^2}.
\]

We now have a complete correspondence:

\[
\text{coefficients} \quad \longleftrightarrow \quad \text{vectors in } H.
\]

\begin{codebox}
\textbf{Notation \texorpdfstring{$\leftrightarrow$}{<->} pygeoinf (analysis/synthesis).}
Fix a concrete \texttt{HilbertSpace} instance \texttt{H}. Its finite working dimension is
\texttt{H.dim} (this is where ``$N$'' enters in \textit{pygeoinf}). Then:
\begin{itemize}
  \item The \emph{primary coefficient route} $\Pi:H\to\ell^2$ is exposed as
  \texttt{H.coordinate\_projection} (a \texttt{LinearOperator} wrapping \texttt{to\_components}).
  \item The synthesis map $\pi:\ell^2\to H$ is exposed as
  \texttt{H.coordinate\_inclusion} (wrapping \texttt{from\_components}).
  \item The ``inner-product vector'' route uses the adjoint
  $\pi^{*_{\ell^2}}:H\to\ell^2$, available as \texttt{H.coordinate\_inclusion.adjoint}.
\end{itemize}
In short, the coefficient vector is \texttt{c = H.coordinate\_projection(m)} (this is $\Pi m$),
and reconstruction is \texttt{m = H.coordinate\_inclusion(c)} (this is $\pi\mathbf{c}$).
\end{codebox}

\section{The Gram Matrix: The Metric Hidden in the Basis}

Up to now we have only discussed linear representation.
We now ask a deeper question:

\medskip
\emph{How is the inner product of $H$ expressed in terms of coefficients?}
\medskip

Let $\mathbf{c} = \Pi m$ and $\mathbf{d} = \Pi n$.
Then
\[
(m,n)_H
=
(\pi \mathbf{c}, \pi \mathbf{d})_H.
\]

If the basis were orthonormal, this would equal $(\mathbf{c},\mathbf{d})_{\ell^2}$.
But for a general Riesz basis this is false.

\subsection{The Induced Coefficient Inner Product}

Define
\begin{equation}
(\mathbf{c},\mathbf{d})_\pi := (\pi \mathbf{c},\pi \mathbf{d})_H .
\end{equation}

This turns $\ell^2$ into a Hilbert space in which
\[
\pi : (\ell^2,(\cdot,\cdot)_\pi)
\to
(H,(\cdot,\cdot)_H)
\]
is an isometry.

Thus the geometry of $H$ is transferred to coefficient space.

\subsection{The Gram Operator}

We now express this induced inner product relative to the
\emph{standard} $\ell^2$ inner product
\[
(\mathbf{c},\mathbf{d})_{\ell^2} = \sum_j [\mathbf{c}]_j [\mathbf{d}]_j.
\]

There exists a unique bounded, self-adjoint, strictly positive operator
\[
\mathbf{M}_\phi : \ell^2 \to \ell^2
\]
such that
\begin{equation}
(\mathbf{c},\mathbf{d})_\pi = (\mathbf{M}_\phi \mathbf{c}, \mathbf{d})_{\ell^2}.
\end{equation}

A direct calculation shows
\begin{equation}
[\mathbf{M}_\phi]_{jk} = (\phi_k,\phi_j)_H,
\qquad
\mathbf{M}_\phi = \pi^{*_{\ell^2}} \pi.
\end{equation}

where $*_{\ell^2}$ denotes adjoint w.r.t. the $\ell^2$ inner product.  The matrix $\mathbf{M}_\phi$ measures how the basis vectors overlap.
It is the \emph{metric tensor} of the representation.

\subsection{Why the Gram Matrix Matters}

The Riesz basis bounds imply
\[
A \|\mathbf{c}\|_{\ell^2}^2
\le
(\mathbf{M}_\phi \mathbf{c},\mathbf{c})_{\ell^2}
\le
B \|\mathbf{c}\|_{\ell^2}^2,
\]
so $\mathbf{M}_\phi$ is symmetric positive definite and invertible.

Moreover, the analysis operator can be computed via
\begin{equation}
\boxed{
\Pi = \mathbf{M}_\phi^{-1} \pi^{*_{\ell^2}}.
}
\end{equation}

This identity is fundamental:

\begin{quote}
\emph{We do not need the dual basis explicitly.
The Gram matrix contains all the necessary information.}
\end{quote}

\begin{codebox}
\textbf{Computing coefficients via the Gram/mass operator in code.}
When you know the synthesis map $\pi$ and can form its Hilbert adjoint
$\pi^{*_{\ell^2}}$, you can recover the analysis operator via
$\Pi = \mathbf{M}_\phi^{-1}\pi^{*_{\ell^2}}$ with
$\mathbf{M}_\phi = \pi^{*_{\ell^2}}\pi$.

In \textit{pygeoinf}, with \texttt{H: HilbertSpace} and \texttt{N = H.dim}:
\begin{itemize}
  \item \texttt{pi = H.coordinate\_inclusion} is $\pi: \mathbb{R}^N\to H$,
  \item \texttt{pi\_star = pi.adjoint} is $\pi^{*_{\ell^2}}: H\to\mathbb{R}^N$,
  \item \texttt{M = pi\_star @ pi} is $\mathbf{M}_{\phi}:\mathbb{R}^N\to\mathbb{R}^N$.
\end{itemize}
Then for a vector $m\in H$:
\begin{enumerate}
  \item Form \texttt{b = pi\_star(m)} (this is the $\pi^{*_{\ell^2}}m$ ``inner-product vector''),
  \item Solve \texttt{M c = b} for \texttt{c} (this is $\Pi m$).
\end{enumerate}
Practical note: use \texttt{.adjoint} for Hilbert adjoints; \texttt{.dual} is the Banach dual ($\pi'$) and differs unless you are in a standard Euclidean setting.
\end{codebox}

\section{Duality and the Riesz Map in Coordinates}

The dual space $H'$ contains functionals.
The Riesz isomorphism
\[
\Riesz_H : H' \to H
\]
identifies each functional with its representing vector.

Its inverse
\[
\Riesz_H^{-1} : H \to H'
\]
turns a vector into its corresponding functional.

\subsection{What Happens in Coordinates}
\label{ssec:riesz-in-coordinates}

Suppose $m \in H$ has coefficient vector $\mathbf{c} = \Pi m \in \ell^2$.
Equivalently, $m = \pi \mathbf{c}$.

We would like to understand the \emph{coordinate representation} of the
functional $\Riesz_H^{-1} m \in H'$.  Concretely, we ask:

\medskip
\emph{If we turn $m$ into a functional via $\Riesz_H^{-1}$, and then extract
its coefficients in coefficient space, what do we get?}
\medskip

Recall:
\begin{itemize}
\item The inverse Riesz map $\Riesz_H^{-1}:H\to H'$ is characterized by
\begin{equation}
\label{eq:riesz-inv-def}
\langle \Riesz_H^{-1} u, v\rangle = (u,v)_H
\qquad \forall\, u,v\in H.
\end{equation}

\item The Banach adjoint (dual operator) of synthesis $\pi:\ell^2\to H$ is
\begin{equation}
\label{eq:pi-prime-def}
\pi' : H' \to (\ell^2)', \qquad
(\pi' \ell)(\mathbf{d}) := \langle \ell, \pi \mathbf{d}\rangle
\qquad \forall\,\ell\in H',\ \mathbf{d}\in \ell^2.
\end{equation}

\item The coefficient-space Riesz map $\Riesz_{\ell^2}:(\ell^2)'\to \ell^2$
(for the \emph{standard} $\ell^2$ inner product) is characterized by
\begin{equation}
\label{eq:riesz-l2-def}
(\Riesz_{\ell^2}\eta,\mathbf{d})_{\ell^2} = \eta(\mathbf{d})
\qquad \forall\,\eta\in(\ell^2)',\ \mathbf{d}\in\ell^2.
\end{equation}

\item The Gram/mass operator $\mathbf{M}_\phi:\ell^2\to\ell^2$ is defined by
\begin{equation}
\label{eq:gram-def-here}
(\pi \mathbf{c},\pi \mathbf{d})_H = (\mathbf{M}_\phi \mathbf{c},\mathbf{d})_{\ell^2}
\qquad \forall\,\mathbf{c},\mathbf{d}\in\ell^2.
\end{equation}
Equivalently, $\mathbf{M}_\phi=\pi^{*_{\ell^2}}\pi$.
\end{itemize}

\paragraph{Derivation of the coordinate identity.}
Let $m\in H$ be arbitrary and set $\mathbf{c}:=\Pi m$ so that $m=\pi\mathbf{c}$.
Define the coefficient functional
\[
\eta := \pi'(\Riesz_H^{-1} m)\in(\ell^2)'.
\]
We compute $\eta(\mathbf{d})$ for an arbitrary coefficient vector
$\mathbf{d}\in\ell^2$:
\begin{align}
\eta(\mathbf{d})
&= \bigl(\pi'(\Riesz_H^{-1} m)\bigr)(\mathbf{d})
&& \text{(definition of $\eta$)} \nonumber\\
&= \left\langle \Riesz_H^{-1} m,\ \pi\mathbf{d}\right\rangle
&& \text{(definition of $\pi'$)} \label{eq:coord-deriv-1}\\
&= (m,\pi\mathbf{d})_H
&& \text{(Riesz characterization \eqref{eq:riesz-inv-def})} \label{eq:coord-deriv-2}\\
&= (\pi\mathbf{c},\pi\mathbf{d})_H
&& \text{(since $m=\pi\mathbf{c}$)} \label{eq:coord-deriv-3}\\
&= (\mathbf{M}_\phi\mathbf{c},\mathbf{d})_{\ell^2}
&& \text{(Gram identity \eqref{eq:gram-def-here}).} \label{eq:coord-deriv-4}
\end{align}
Now compare \eqref{eq:coord-deriv-4} with the defining property of the
coefficient-space Riesz map \eqref{eq:riesz-l2-def}.
Since \eqref{eq:riesz-l2-def} says that $\Riesz_{\ell^2}\eta$ is the unique
vector in $\ell^2$ satisfying
\[
(\Riesz_{\ell^2}\eta,\mathbf{d})_{\ell^2}=\eta(\mathbf{d})\quad\forall\,\mathbf{d},
\]
and we have just shown that $\eta(\mathbf{d})=(\mathbf{M}_\phi\mathbf{c},\mathbf{d})_{\ell^2}$
for all $\mathbf{d}$, it follows that
\[
\Riesz_{\ell^2}\eta = \mathbf{M}_\phi \mathbf{c}.
\]
Substituting back $\eta=\pi'(\Riesz_H^{-1} m)$ and $\mathbf{c}=\Pi m$ yields the
operator identity
\begin{equation}
\label{eq:riesz-coord-identity}
\boxed{
\Riesz_{\ell^2}\,\pi'\,\Riesz_H^{-1}
=
\mathbf{M}_\phi\,\Pi.
}
\end{equation}

\paragraph{Interpretation: why this means $\mathbf{M}_\phi$ is the coordinate representation of $\Riesz_H^{-1}$.}
Equation \eqref{eq:riesz-coord-identity} is an equality of maps $H\to \ell^2$.
Read it from right to left and from left to right:

\begin{itemize}
\item The map $\Pi:H\to\ell^2$ takes a vector $m$ and returns its \emph{primal}
coefficient vector $\mathbf{c}=\Pi m$.

\item Multiplication by $\mathbf{M}_\phi$ converts this primal coefficient
vector into $\mathbf{M}_\phi\mathbf{c}\in\ell^2$.

\item The composite $\Riesz_{\ell^2}\,\pi'\,\Riesz_H^{-1}$ does the same job
but via the dual space: it first turns $m$ into its functional
$\Riesz_H^{-1}m\in H'$, then extracts the corresponding coefficient functional
$\pi'(\Riesz_H^{-1}m)\in(\ell^2)'$, and finally identifies that functional with
a coefficient vector in $\ell^2$ using $\Riesz_{\ell^2}$.

\end{itemize}

Thus, after choosing coordinates via $\Pi$ (on $H$) and via
$\Riesz_{\ell^2}\circ \pi'$ (on $H'$), the inverse Riesz map
$\Riesz_H^{-1}:H\to H'$ is represented by the operator $\mathbf{M}_\phi$ on the
coefficient space.  In this precise sense,

\begin{center}
\textbf{$\mathbf{M}_\phi$ is the coordinate (matrix/operator) representation of $\Riesz_H^{-1}$.}
\end{center}

Equivalently, one may solve \eqref{eq:riesz-coord-identity} for $\Riesz_H^{-1}$:
\begin{equation}
\label{eq:riesz-inv-from-gram}
\Riesz_H^{-1}
=
(\pi')^{-1}\,\Riesz_{\ell^2}^{-1}\,\mathbf{M}_\phi\,\Pi,
\end{equation}
which makes explicit that $\mathbf{M}_\phi$ is the ``middle'' coordinate map
that carries primal coefficients to dual coefficients.


\begin{figure}[t]
\centering
\begin{tikzcd}[column sep=large,row sep=large]
H \arrow[rr,"\Pi_N"] \arrow[rd,swap,"\Pi"] \arrow[ddd, swap, bend right=25, "\Riesz_H^{-1}"] & & \mathbb{R}^N \arrow[ll, bend right=25, swap, "\pi_N"] \arrow[ddd, bend left=25, "\mathbf{M}_{\phi, N}"]\\
& \ell^2 \arrow[lu, bend right=25, swap, "\pi"]  \arrow[d, bend right=25, swap, "\Riesz_{\ell^2}^{-1} \mathbf{M}_{\phi}"] & \\
& (\ell^2)' \arrow[dl, swap, "\Pi'"] \arrow[u, swap, "\Riesz_{\ell^2} \mathbf{M}_{\phi}^{-1}"]& \\
H' \arrow[rr, swap, "\pi_N'"] \arrow[ur, swap, bend right=25, "\pi'"] \arrow[uuu,"\Riesz_H"] & &
(\mathbb{R}^N)' \arrow[ll, bend left=25, "\Pi_N'"] \arrow[uuu, "\mathbf{M}_{\phi, N}^{-1}"]
\end{tikzcd}
\caption{Commutative diagram of primal and dual coefficient maps in \textit{pygeoinf}.
The analysis operator $\Pi: H \to \ell^2$ extracts primal coefficients, while its inverse (synthesis) $\pi: \ell^2 \to H$ reconstructs vectors .
The Riesz isomorphism $\mathcal{R}_H: H' \to H$ identifies dual functionals with primal representatives; its inverse $\mathcal{R}_H^{-1}: H \to H'$ maps a vector to its unique functional form $\langle \mathcal{R}_H^{-1} u, v \rangle = (u, v)_H$ .
In the coefficient space, the Gram matrix $\mathbf{M}_{\phi}$ acts as the coordinate representation of the inverse Riesz map, transforming primal coefficients into dual coefficients .
The diagram commutes such that extracting dual coefficients directly from a functional is equivalent to mapping primal coefficients through the Gram matrix: $\Riesz_{\ell^2} \circ \pi' \circ \mathcal{R}_H^{-1} = \mathbf{M}_{\phi} \circ \Pi$ (identifying $(\ell^2)' \cong \ell^2$ via $\Riesz_{\ell^2}$).}
\label{fig:primal-dual-coeff}
\end{figure}

\section{Finite-Dimensional Approximation}

In practice we never work with infinite coefficient sequences.
Instead, we truncate the basis to the first $N$ elements and work in the
finite-dimensional subspace
\begin{equation}
H_N := \mathrm{span}\{\phi_1,\dots,\phi_N\} \subset H.
\end{equation}

Everything in the infinite-dimensional theory now reduces to linear algebra.

\subsection{Finite Synthesis and Analysis}

Define the truncated synthesis operator
\begin{equation}
\pi_N : \mathbb{R}^N \to H_N,
\qquad
\pi_N(\mathbf{c}) := \sum_{j=1}^N [\mathbf{c}]_j \phi_j .
\end{equation}

Since $\{\phi_1,\dots,\phi_N\}$ is linearly independent,
$\pi_N$ is a linear isomorphism.

Its inverse is the truncated analysis operator
\begin{equation}
\Pi_N : H_N \to \mathbb{R}^N,
\qquad
\Pi_N(m) = \mathbf{c} \quad \text{such that } m=\pi_N(\mathbf{c}).
\end{equation}

Thus
\begin{equation}
\pi_N \circ \Pi_N = \mathrm{Id}_{H_N},
\qquad
\Pi_N \circ \pi_N = \mathrm{Id}_{\mathbb{R}^N}.
\end{equation}

\subsection{The Finite Gram Matrix}

The geometry of $H_N$ is encoded by the
$N\times N$ Gram (mass) matrix
\begin{equation}
[\mathbf{M}_{\phi,N}]_{jk}
:=
(\phi_k,\phi_j)_H,
\qquad 1 \le j,k \le N.
\end{equation}

A direct calculation shows
\begin{equation}
(\pi_N \mathbf{c},\pi_N \mathbf{d})_H
=
(\mathbf{c}, \mathbf{M}_{\phi,N} \mathbf{d})_{\mathbb{R}^N},
\end{equation}
where $(\mathbf{c},\mathbf{d})_{\mathbb{R}^N}=\mathbf{c}^\top \mathbf{d}$ is the standard Euclidean inner product.

Thus $M_{\phi,N}$ represents the induced inner product
\[
(\mathbf{c},\mathbf{d})_\pi := (\pi_N \mathbf{c},\pi_N \mathbf{d})_H.
\]

\subsection{Properties of the Finite Gram Matrix}

The matrix $\mathbf{M}_{\phi,N}$ is:

\begin{itemize}
\item symmetric,
\item strictly positive definite,
\item invertible.
\end{itemize}

Indeed,
\[
\mathbf{c}^\top \mathbf{M}_{\phi,N} \mathbf{c}
=
\|\pi_N \mathbf{c}\|_H^2,
\]
and this vanishes only if $\mathbf{c}=0$.

\subsection{Finite Analysis via the Gram Matrix}

Taking adjoints relative to the standard Euclidean inner product,
\[
\pi_N^* : H_N \to \mathbb{R}^N,
\qquad
(\pi_N^* m)_j = (m,\phi_j)_H.
\]

The truncated analysis operator satisfies
\begin{equation}
\boxed{
\Pi_N = \mathbf{M}_{\phi,N}^{-1} \pi_N^*.
}
\end{equation}

Thus coefficients can be computed without explicitly constructing
the dual basis:
\begin{enumerate}
\item Compute inner products with $\phi_j$.
\item Solve a linear system with matrix $\mathbf{M}_{\phi,N}$.
\end{enumerate}

\subsection{Finite Riesz Map in Coordinates}

The Riesz isomorphism restricted to $H_N$ satisfies
\[
\langle \Riesz_H^{-1} m , v \rangle
=
(m,v)_H,
\qquad
m,v \in H_N.
\]

In coefficient form, let $\mathbf{c}=\Pi_N m$.
Then the dual coefficients of $\Riesz_H^{-1} m$ are
\begin{equation}
\boxed{
\Riesz_{\mathbb{R}^N}\,\pi_N'\,\Riesz_H^{-1}
=
\mathbf{M}_{\phi,N}\,\Pi_N.
}
\end{equation}

In purely matrix terms:

\begin{equation}
\text{dual coefficients}
=
\mathbf{M}_{\phi,N}\,(\text{primal coefficients}).
\end{equation}

Thus the finite Gram matrix is the matrix representation
of the inverse Riesz map on $H_N$.

\subsection{Summary: What Changes in Finite Dimensions}

In infinite dimensions:
\begin{itemize}
\item $\ell^2$ and $H$ are infinite Hilbert spaces.
\item $\mathbf{M}_\phi$ is a bounded positive operator.
\item Riesz maps are abstract isomorphisms.
\end{itemize}

In finite dimensions:
\begin{itemize}
\item All spaces are Euclidean.
\item All operators are matrices.
\item The Riesz map becomes multiplication by a symmetric positive definite matrix.
\item The Gram matrix \emph{is} that matrix.
\end{itemize}

Thus the finite-dimensional setting is not an approximation of the structure ---
it is the exact same structure realized in linear algebra form.

\begin{whybox}
This is the key theory\,$\leftrightarrow$\,computation bridge: the abstract operators
($\pi$, $\Pi$, adjoints, and Riesz maps) do not disappear when you discretize.
Choosing a concrete \texttt{HilbertSpace} (and therefore a specific \texttt{dim}) \emph{is} the discretization.
Once that choice is made, \textit{pygeoinf} computes adjoints using the declared inner products,
so the same identities (e.g.\ $\Pi = \mathbf{M}^{-1}\pi^{*_{\ell^2}}$) remain valid and numerically meaningful.
As you increase \texttt{dim} (or refine the underlying space), you are moving along a controlled family of finite problems that approximate the intended infinite-dimensional model.
\end{whybox}
\section{Concrete spaces in pygeoinf}
\label{sec:concrete-spaces}

This section introduces the built-in Hilbert space implementations and
explains how to create your own.

\subsection{\texttt{EuclideanSpace}}
\label{ssec:euclidean-space}

$\mathbb{R}^n$ with the standard inner product $(u, v) = \sum_i u_i v_i =
\mathbf{u}^\top \mathbf{v}$. The Riesz map is the identity: $\Riesz = \mathrm{Id}$.
The basis $\{\phi_j\} = \{\mathbf{e}_j\}$ is the standard ONB, so $\mathbf{M}_\phi
= \mathbf{I}$ and $\Pi = \pi^{*_{\ell^2}} = \mathrm{Id}$.

When you use \texttt{EuclideanSpace(n)}, adjacent operator transposes
\emph{are} adjoints, and you recover the classical formulas of linear
algebra. This is the appropriate choice when variables represent
coefficients in a pre-orthonormalised basis (e.g.\ after a change of
variables $\tilde{m} = \mathbf{L} m$ where $\mathbf{L}$ is a Cholesky
factor of $\mathbf{M}_\phi$).

\subsection{\texttt{MassWeightedHilbertSpace}}
\label{ssec:mass-weighted}

Given a Hilbert space $H$ and a bounded symmetric positive operator
$\mathbf{W} : H \to H'$ (the ``weight operator''), the
\emph{mass-weighted space} $H_{\mathbf{W}}$ is the same underlying vector
space but with inner product
\begin{equation}
  \label{eq:mass-weighted-ip}
  (u, v)_{H_{\mathbf{W}}} = \langle \mathbf{W} u,\, v \rangle.
\end{equation}
The Riesz maps become $\Riesz_{\mathbf{W}} = \Riesz \circ \mathbf{W}$ and
$\Riesz_{\mathbf{W}}^{-1} = \mathbf{W}^{-1} \circ \Riesz^{-1}$.

\paragraph{Motivating use case: discretized $L^2$.}
Using the non-uniform grid inner product from
\cref{sec:why-spaces-matter}, we set $H = \mathbb{R}^N$,
$H_{\mathbf{W}} = \mathbb{R}^N$ with weight $\mathbf{W} = \mathbf{M} =
\mathrm{diag}(\Delta x_j)$. Then $(u, v)_{H_{\mathbf{W}}} = \sum_j
u_j v_j \Delta x_j \approx \int_0^1 u(x) v(x)\,\mathrm{d}x$.

\subsection{\texttt{HilbertSpaceDirectSum}}
\label{ssec:direct-sum-space}

Given Hilbert spaces $H_1$ and $H_2$, their \emph{Hilbert direct sum}
$H_1 \oplus H_2$ consists of pairs $(u_1, u_2)$ with
\begin{equation}
  \label{eq:direct-sum-ip}
  \bigl((u_1, u_2),\, (v_1, v_2)\bigr)_{H_1 \oplus H_2}
  = (u_1, v_1)_{H_1} + (u_2, v_2)_{H_2}.
\end{equation}
The Riesz maps and analysis/synthesis operators decompose blockwise.
This is the natural structure for models that are concatenations of
separately-structured fields (e.g.\ $P$-wave and $S$-wave velocity
perturbations).

\subsection{Implementing your own \texttt{HilbertSpace}}
\label{ssec:custom-space}

To encode a non-standard model space, subclass \texttt{HilbertSpace}.
The abstract interface requires exactly eight methods:

\begin{codebox}
Subclass \texttt{HilbertSpace} and implement:
\begin{itemize}
  \item \texttt{dim} --- integer dimension $N$ of the space.
  \item \texttt{zero()} --- return the zero element.
  \item \texttt{basis\_vector(j)} --- return $\phi_j$ for $j = 0, \ldots, N-1$.
  \item \texttt{inner\_product(x1, x2)} --- compute $(x_1, x_2)_H$.
  \item \texttt{to\_dual(x)} --- compute $\Riesz^{-1}(x) \in H'$.
  \item \texttt{from\_dual(xp)} --- compute $\Riesz(\ell) \in H$ for $\ell \in H'$.
  \item \texttt{to\_components(x)} --- compute $\Pi(x) \in \mathbb{R}^N$.
  \item \texttt{from\_components(c)} --- compute $\pi(\mathbf{c}) \in H$.
\end{itemize}
All other methods (\texttt{norm}, \texttt{squared\_norm}, \texttt{distance},
\texttt{riesz}, \texttt{inverse\_riesz}, \texttt{coordinate\_projection},
\texttt{coordinate\_inclusion}, \ldots) are derived automatically from the above.
\end{codebox}

\begin{examplebox}
\textbf{Example (quadrature-weighted $L^2$ via \texttt{MassWeightedHilbertSpace}).}
Suppose you discretize $[0,1]$ on a non-uniform grid with cell widths
$\Delta x_j$ and represent a function by its sampled values
$\mathbf{u}=(u_j)\in\mathbb{R}^N$.
The quadrature-weighted inner product
\[
(\mathbf{u},\mathbf{v}) \approx \sum_{j=1}^N u_j v_j\,\Delta x_j
\]
is obtained by taking an underlying Euclidean space $X=\mathbb{R}^N$ and a diagonal
mass operator $\mathbf{M}=\mathrm{diag}(\Delta x_j)$, then forming the mass-weighted space
$X_{\mathbf{M}}$.

In \textit{pygeoinf} this is the canonical pattern:
\begin{itemize}
  \item \texttt{X = EuclideanSpace(N)}
  \item \texttt{dx = np.asarray([\ldots])  \# quadrature weights (e.g. cell widths)}
  \item \texttt{M = LinearOperator.from\_matrix(X, X, np.diag(dx))}
  \item \texttt{Minv = LinearOperator.from\_matrix(X, X, np.diag(1.0/dx))}
  \item \texttt{H = MassWeightedHilbertSpace(X, M, Minv)}
\end{itemize}
Then \texttt{H.inner\_product(u, v)} computes $\langle \mathbf{M}\mathbf{u},\mathbf{v}\rangle$
and the Riesz maps become $\Riesz_H = \Riesz_X\circ\mathbf{M}$ and
$\Riesz_H^{-1}=\mathbf{M}^{-1}\circ\Riesz_X^{-1}$, exactly matching
\eqref{eq:mass-weighted-ip}.
\end{examplebox}



% ============================================================================
\chapter{Linear Forms}
\label{ch:linear-forms}
% ============================================================================

\section{Motivation: measurements as linear functionals}
\label{sec:forms-motivation}

Many inverse problems can be written (after linearization, if needed) as
\[
  \tilde{d}_i = \ell_i(m) + \varepsilon_i, \qquad i=1,\ldots,M,
\]
where $m\in\modelspace$ is the model, each $\ell_i\in\modelspace'$ is a
continuous linear functional, and $\varepsilon_i$ is measurement noise.
The data vector is therefore best understood as a \emph{stack of linear forms}
acting on the model.

\begin{examplebox}
\textbf{Example (a localized average).}
Let $\modelspace=L^2(\Omega)$ and let $w\in L^2(\Omega)$ be a fixed weighting
function describing an instrument footprint or averaging kernel.
Then
\[
  \ell(m) := \int_\Omega w(x)\,m(x)\,\mathrm{d}x
\]
is a continuous linear form.
In a discretization, the only question is how $m$ is represented and which
inner product defines the geometry of $\modelspace$.
\end{examplebox}

\begin{whybox}
Every single datum in your data vector is the result of applying a
linear functional to the model. Understanding linear forms is
understanding what your data actually \emph{are} mathematically.
\end{whybox}

\section{Definition}
\label{sec:form-definition}

Let $\modelspace$ be a (real) Hilbert space and let $\modelspace'$ denote its
continuous dual.
A \emph{linear form} is a continuous linear map
\[
  \ell : \modelspace \to \mathbb{R}.
\]
We write the duality pairing as
\[
  \pair{\ell}{m} := \ell(m), \qquad \ell\in\modelspace',\; m\in\modelspace.
\]

In \textit{pygeoinf}, a linear form is represented by the class
\texttt{LinearForm(domain, \ldots)}, where \texttt{domain} is the Hilbert space
on which the form acts.

\section{Representation in coordinates}
\label{sec:form-coordinates}

The library uses the domain's coordinate maps
$\Pi:\modelspace\to\mathbb{R}^N$ (\texttt{to\_components}) and
$\pi:\mathbb{R}^N\to\modelspace$ (\texttt{from\_components}) to provide a
computational coordinate system.
This is a \emph{chosen representation}: it need not be orthonormal.

Let $\phi_j := \pi(e_j)$ denote the corresponding computational basis vectors
(in code: \texttt{domain.basis\_vector(j)}), where $e_j$ are the standard unit
vectors in $\mathbb{R}^N$.
The component vector of a linear form $\ell$ is defined as
\[
  \mathbf{c}_\ell := (\ell(\phi_0),\ldots,\ell(\phi_{N-1}))^\top \in \mathbb{R}^N.
\]
Then for any $m\in\modelspace$ with coefficient vector $\mathbf{m}:=\Pi m$,
linearity gives
\[
  \ell(m) = \sum_{j=0}^{N-1} \ell(\phi_j)\,\mathbf{m}_j
  = \mathbf{c}_\ell^\top\,\mathbf{m}.
\]
This is exactly the evaluation rule used by \texttt{LinearForm}:
\texttt{l(m) = dot(l.components, domain.to\_components(m))}.

\subsection*{Two equivalent construction paths in the API}

\begin{codebox}
\textbf{(1) Provide components directly.}\\
\texttt{l = LinearForm(domain, components=c)}\\
Here \texttt{c} is a 1D array of length \texttt{domain.dim} representing
$\mathbf{c}_\ell$.

\medskip
\textbf{(2) Provide a callable mapping and let pygeoinf compute components.}\\
\texttt{l = LinearForm(domain, mapping=f, parallel=True, n\_jobs=-1)}\\
The library computes \texttt{l.components[i]} by evaluating \texttt{f} on the
computational basis vectors $\phi_i=\texttt{domain.basis\_vector(i)}$
(equivalently $\phi_i=\texttt{domain.from\_components}(e_i)$):
$\mathbf{c}_{\ell,i} = f(\phi_i)$.
When \texttt{parallel=True}, this is done with a joblib backend; \texttt{n\_jobs}
controls the number of workers.
\end{codebox}

\begin{pitfallbox}
The \texttt{mapping} construction is convenient, but it assumes the callable
really is linear (and deterministic) on the domain.
If your \texttt{mapping} includes hidden nonlinearities (thresholding, clipping,
normalization, stateful caches), the computed components will be inconsistent
with the linear algebra and subsequent results may be meaningless.
\end{pitfallbox}

\section{The Riesz representative (gradient)}
\label{sec:form-riesz-rep}

Because $\modelspace$ is a Hilbert space, every $\ell\in\modelspace'$ has a
unique \emph{Riesz representative} $g_\ell\in\modelspace$ such that
\[
  \ell(m) = \pair{\ell}{m} = (g_\ell, m)_{\modelspace}\qquad\text{for all } m\in\modelspace.
\]
Equivalently, $g_\ell = \Riesz_{\modelspace}(\ell)$, where
$\Riesz_{\modelspace}:\modelspace'\to\modelspace$ is the Riesz isomorphism
listed in the notation table (\texttt{from\_dual}).

In \textit{pygeoinf}, this is the mechanism that injects \emph{geometry} into
linear forms: the component vector $\mathbf{c}_\ell$ encodes the action of
$\ell$ on the chosen coordinate system, while the inner product on
$\modelspace$ determines the corresponding gradient vector.
Concretely,
\[
  \nabla\ell \equiv g_\ell = \texttt{domain.from\_dual(l)}.
\]
The Hessian of any linear form is the zero operator.

\begin{remark}
If the coordinate system is not orthonormal, then the coefficient vector
$\Pi g_\ell$ is generally \emph{not} equal to $\mathbf{c}_\ell$.
They are related by the Gram/mass matrix of the chosen basis.
This is why \texttt{from\_dual} (Riesz) must be used to obtain the correct
gradient in the domain geometry.
\end{remark}

\section{Interoperability: viewing a linear form as an operator}
\label{sec:form-as-operator}

Some parts of the library (and many external numerical tools) are written in
terms of operators rather than functionals.
For interoperability, \texttt{LinearForm} exposes an adapter
\texttt{as\_linear\_operator} that views $\ell$ as an operator
$\modelspace\to\mathbb{R}$ by returning a \texttt{LinearOperator} with codomain
\texttt{EuclideanSpace(1)}.

This should be understood purely as an interface bridge, not as a conceptual
re-framing of linear forms.
In particular, we do \emph{not} treat a linear form as a ``rank-1 operator'' in
the theory: its natural home is the dual space $\modelspace'$.

\subsection*{Implementing your own linear form (practical guidance)}

Most users should not subclass anything: instantiate \texttt{LinearForm} using
one of the two construction paths above.
To ensure correctness, keep the following invariants in mind:
\begin{itemize}
  \item \textbf{Component meaning.} The stored array \texttt{l.components} is
    $\mathbf{c}_\ell$ with entries $\mathbf{c}_{\ell,i}=\ell(\phi_i)$ for the
    domain's chosen basis vectors $\phi_i=\texttt{domain.basis\_vector(i)}$.
  \item \textbf{Evaluation rule.} Evaluation uses the coordinate projection:
    $\ell(m)=\mathbf{c}_\ell^\top\,\Pi m$.
    If you change how \texttt{to\_components/from\_components} represent the
    domain, you also change what \texttt{components} mean.
  \item \textbf{Geometry lives in the Riesz map.} Use
    \texttt{domain.from\_dual(l)} to obtain the gradient/Riesz representative.
\end{itemize}

Convenience methods commonly used in implementations and experiments:
\begin{itemize}
  \item Arithmetic: \texttt{l1 + l2}, \texttt{l1 - l2}, \texttt{a * l},
    \texttt{l / a}, and in-place \texttt{l *= a}, \texttt{l += other}.
  \item \texttt{l.copy()} for a deep copy.
  \item \texttt{l.as\_linear\_operator} to plug a form into operator-based code.
  \item \texttt{LinearForm.from\_linear\_operator(op)} to wrap an operator
    \texttt{op: domain -> EuclideanSpace(1)} as a linear form.
\end{itemize}

\begin{codebox}
\texttt{LinearForm(domain, components=...)} stores $\mathbf{c}_\ell$ and evaluates
via \texttt{dot(components, domain.to\_components(x))}.\\
\texttt{LinearForm(domain, mapping=..., parallel=..., n\_jobs=...)} computes
components by evaluating the mapping on \texttt{domain.basis\_vector(i)}.

\medskip
Useful methods/properties: \texttt{\_\_call\_\_(x)}, \texttt{components},
\texttt{copy()}, \texttt{as\_linear\_operator}, \texttt{from\_linear\_operator(...)}.
\end{codebox}


% ============================================================================
\chapter{Linear Operators}
\label{ch:linear-operators}
% ============================================================================

\section{Motivation: the forward operator}
\label{sec:operator-motivation}
In inverse problems we distinguish a \emph{model space} $\modelspace$ and a
\emph{data space} $\dataspace$. The forward physics is encoded as a linear map
$G: \modelspace \to \dataspace$ taking a model $m$ to the data $d = G(m)$ it
predicts.
Throughout this chapter we use the evaluation convention $G(x)$ (rather than
juxtaposition $Gx$) to emphasize that operators are functions, and to match the
callable operator objects in the \texttt{pygeoinf} API.

\begin{whybox}
The forward operator $G$ is the mathematical encoding of your physics.
Getting its adjoint right is critical for every inversion algorithm ---
and the adjoint depends on the inner products of both spaces, not just
the matrix transpose.
\end{whybox}

\section{Definition and properties}
\label{sec:operator-definition}

\paragraph{Definition (bounded linear operator).}
Let $\modelspace$ and $\dataspace$ be normed vector spaces (in particular,
Hilbert spaces in the applications in this manual). A map
$G: \modelspace \to \dataspace$ is \emph{linear} if
\begin{equation}
  G(\alpha m_1 + \beta m_2) = \alpha\,G(m_1) + \beta\,G(m_2)
  \qquad \forall\; m_1,m_2\in\modelspace,\; \alpha,\beta\in\mathbb{R}.
\end{equation}
It is \emph{bounded} if there exists a constant $C < \infty$ such that
\begin{equation}
  \label{eq:bounded-linear-operator}
  \|G(m)\|_{\dataspace} \le C\,\|m\|_{\modelspace}
  \qquad \forall\; m\in\modelspace.
\end{equation}
The smallest such constant is the \emph{operator norm}
$\|G\| = \sup_{\|m\|_{\modelspace}=1} \|G(m)\|_{\dataspace}$.
In Hilbert spaces, boundedness is equivalent to continuity, and it is the
assumption that makes adjoints and discretisations behave well.

\begin{codebox}
Notation \(\leftrightarrow\) \texttt{pygeoinf} API (linear operators):
\[
  G: \modelspace \to \dataspace
  \;\equiv\;
  \texttt{LinearOperator(domain, codomain, mapping)}
  \quad \text{and evaluate via } \texttt{G(m)}.
\]
\[
  G'\;\text{(Banach dual)} \;\equiv\; \texttt{G.dual},
  \qquad
  G^*\;\text{(Hilbert adjoint)} \;\equiv\; \texttt{G.adjoint}.
\]
\[
  \Riesz_{\modelspace} \equiv \texttt{domain.from\_dual},\quad
  \Riesz_{\modelspace}^{-1} \equiv \texttt{domain.to\_dual},
  \qquad
  \Riesz_{\dataspace} \equiv \texttt{codomain.from\_dual},\quad
  \Riesz_{\dataspace}^{-1} \equiv \texttt{codomain.to\_dual}.
\]
\end{codebox}

\section{The Banach dual (transpose) $G'$}
\label{sec:banach-dual}

The definition of $G'$ requires only dual pairings; it makes sense for any
bounded linear operator between normed spaces.

\paragraph{Typing.}
Let $\modelspace$ and $\dataspace$ be normed vector spaces and let
$G : \modelspace \to \dataspace$ be a bounded linear operator.
Then the continuous dual spaces $\modelspace'$ and $\dataspace'$ are Banach
spaces (\cref{def:dual-space}), and the dual pairing is denoted by
$\pair{\cdot}{\cdot}$.

\paragraph{Definition (Banach dual / transpose).}
The \emph{Banach dual} (also called the \emph{transpose}) of $G$ is the bounded
linear map
\begin{equation}
  \label{eq:banach-dual-typing}
  G' : \dataspace' \to \modelspace'
\end{equation}
defined by
\begin{equation}
  \label{eq:banach-dual-def}
  \pair{G'(d')}{m} := \pair{d'}{G(m)}
  \qquad \forall\; d'\in\dataspace',\; m\in\modelspace.
\end{equation}
Equivalently, for each $d'\in\dataspace'$, the element $G'(d')\in\modelspace'$
is the functional
\begin{equation}
  (G'(d'))(m) = d'(G(m)).
\end{equation}

\paragraph{Boundedness.}
From \eqref{eq:banach-dual-def} and \eqref{eq:bounded-linear-operator}, for any
$d'\in\dataspace'$ we have
\begin{equation}
  \|G'(d')\|_{\modelspace'}
  = \sup_{\|m\|_{\modelspace}=1} \bigl|\pair{G'(d')}{m}\bigr|
  = \sup_{\|m\|_{\modelspace}=1} \bigl|\pair{d'}{G(m)}\bigr|
  \le \|d'\|_{\dataspace'}\,\|G\|,
\end{equation}
so $G'$ is bounded with $\|G'\|\le\|G\|$.

\paragraph{Why $G'$ is not an adjoint.}
The map $G'$ acts on \emph{dual} elements and produces a dual element.
Even when $\modelspace$ and $\dataspace$ are Hilbert spaces, the Hilbert adjoint
$G^*$ (next section) is a map $\dataspace \to \modelspace$ between the primal
spaces. To pass from $G'$ to $G^*$ we must use the Riesz maps, i.e.\ identify
dual vectors with their Riesz representatives in the primal spaces.

\section{The Hilbert adjoint $G^*$}
\label{sec:hilbert-adjoint}

From here on we assume that both spaces are Hilbert spaces, so that we have
inner products and Riesz maps available.

\paragraph{Typing.}
Let $\modelspace$ and $\dataspace$ be Hilbert spaces with inner products
$\ipM{\cdot}{\cdot}$ and $\ipD{\cdot}{\cdot}$, and let
$G : \modelspace \to \dataspace$ be bounded and linear.

\paragraph{Definition (Hilbert adjoint).}
The \emph{Hilbert adjoint} of $G$ is the unique bounded linear operator
\begin{equation}
  \label{eq:hilbert-adjoint-typing}
  G^* : \dataspace \to \modelspace
\end{equation}
such that
\begin{equation}
  \label{eq:hilbert-adjoint-def}
  \ipD{G(m)}{d} = \ipM{m}{G^*(d)}
  \qquad \forall\; m\in\modelspace,\; d\in\dataspace.
\end{equation}

\paragraph{Existence and construction (Riesz representation).}
Fix $d\in\dataspace$ and define a linear functional on $\modelspace$ by
\begin{equation}
  \label{eq:adjoint-functional}
  \ell_d(m) := \ipD{G(m)}{d}.
\end{equation}
Boundedness of $G$ implies $\ell_d\in\modelspace'$. By the Riesz theorem
(\cref{thm:riesz}) there exists a unique element $u\in\modelspace$ such that
$\ell_d(m)=\ipM{m}{u}$ for all $m$. We define $G^*(d):=u$.

\paragraph{Derivation of the commuting relation.}
Let $\Riesz_{\modelspace}:\modelspace'\to\modelspace$ and
$\Riesz_{\dataspace}:\dataspace'\to\dataspace$ denote the Riesz maps, with
inverses $\Riesz_{\modelspace}^{-1}:\modelspace\to\modelspace'$ and
$\Riesz_{\dataspace}^{-1}:\dataspace\to\dataspace'$.
Using the identity (\cref{sec:dual-space}, \cref{sec:riesz-theorem})
\begin{equation}
  \label{eq:pairing-via-riesz}
  \pair{\ell}{x} = \ip{x}{\Riesz\ell}{X}
  \qquad (\ell\in X',\; x\in X),
\end{equation}
we compute, for arbitrary $m\in\modelspace$ and $d\in\dataspace$,
\begin{align}
  \ipD{G(m)}{d}
  &= \pair{\Riesz_{\dataspace}^{-1} d}{G(m)}
    && \text{by \eqref{eq:pairing-via-riesz} in $X=\dataspace$} \\
  &= \pair{G'\bigl(\Riesz_{\dataspace}^{-1} d\bigr)}{m}
    && \text{by the definition of $G'$ \eqref{eq:banach-dual-def}} \\
  &= \ipM{m}{\Riesz_{\modelspace}\,G'\,\Riesz_{\dataspace}^{-1}(d)}
    && \text{by \eqref{eq:pairing-via-riesz} in $X=\modelspace$.}
\end{align}
Comparing with \eqref{eq:hilbert-adjoint-def} and using uniqueness in the Riesz
representation, we obtain the key relation
\begin{equation}
  \label{eq:adjoint-from-dual}
  G^* = \Riesz_{\modelspace} \circ G' \circ \Riesz_{\dataspace}^{-1}.
\end{equation}

\paragraph{API correspondence.}
In \texttt{pygeoinf}, this is exactly the rule that turns the Banach dual into a
Hilbert adjoint:
\begin{equation}
  \texttt{G.adjoint} \;=\; \texttt{domain.from\_dual}\,\circ\,\texttt{G.dual}\,\circ\,\texttt{codomain.to\_dual}.
\end{equation}

\begin{pitfallbox}
If you are used to thinking of the matrix transpose as the adjoint, beware:
that is only true in Euclidean space with the standard inner product.
For general Hilbert spaces, $G^* \neq G^\top$. The correct relationship is
$G^* = \Riesz_{\modelspace} \circ G' \circ \Riesz_{\dataspace}^{-1}$.
\end{pitfallbox}

\section{Matrix representations}
\label{sec:matrix-representations}
% Standard matrix: A = Π_D G π_N.
% Galerkin matrix.
% Discrete adjoint: A* = M_φ⁻¹ Aᵀ M_b.

\subsection{The standard (component) matrix}
\label{ssec:standard-matrix}

\subsection{The Galerkin matrix}
\label{ssec:galerkin-matrix}

\subsection{The discrete adjoint formula}
\label{ssec:discrete-adjoint}

\section{Operators as stacks of linear forms}
\label{sec:operator-from-forms}
% G = [ℓ₁, ..., ℓ_M]ᵀ with ℓᵢ ∈ M'.

\begin{codebox}
Given linear forms $\ell_i \in \modelspace'$, construct the forward operator via
\texttt{LinearOperator.from\_linear\_forms([$\ell_1, \ldots, \ell_M$])}.
\end{codebox}

\section{Constructing operators in pygeoinf}
\label{sec:constructing-operators}
% From callables, from matrices, from linear forms.
% Providing vs deriving adjoint/dual mappings.

\begin{codebox}
\texttt{LinearOperator(domain, codomain, mapping)} is the primary constructor.
Optionally provide \texttt{dual\_mapping} ($G'$) or \texttt{adjoint\_mapping} ($G^*$);
the library derives the other via the Riesz maps.
\end{codebox}

\section{The full commutative diagram}
\label{sec:operator-diagram}
% Complete tikz-cd diagram: M, D, M', D', ℝⁿ, ℝ^{N_d}, all maps.


% ============================================================================
\chapter{Nonlinear Operators and Forms}
\label{ch:nonlinear}
% ============================================================================

\section{NonLinearOperator}
\label{sec:nonlinear-operator}
% Interface: __call__, jacobian, adjoint_jacobian.

\section{NonLinearForm}
\label{sec:nonlinear-form}
% Scalar-valued nonlinear functions on Hilbert spaces.

\section{Gradients and Hessians}
\label{sec:frechet-derivatives}
% Fréchet derivatives, connection to optimisation.


% %%%%%%%%%%%%%%%%%%%%%%%%%%%%%%%%%%%%%%%%%%%%%%%%%%%%%%%%%%%%%%%%%%%%%%%%%%%%
\part{Geometry and Structure}
\label{part:geometry}
% %%%%%%%%%%%%%%%%%%%%%%%%%%%%%%%%%%%%%%%%%%%%%%%%%%%%%%%%%%%%%%%%%%%%%%%%%%%%

% ============================================================================
\chapter{Subsets and Convex Sets}
\label{ch:subsets}
% ============================================================================

\section{Why convex sets?}
\label{sec:why-convex}
% Prior information is geometric: the model lies in some set.
% Convexity = computational tractability.

\begin{whybox}
Every piece of prior information you have about your model ---
bounds on parameters, energy constraints, smoothness requirements ---
defines a set of admissible models. Convexity is the mathematical
property that makes these constraints computationally tractable.
\end{whybox}

\section{The Subset interface}
\label{sec:subset-interface}
% is_element(), project(), boundary.

\section{Balls and ellipsoids}
\label{sec:balls-ellipsoids}
% Definition, projection formulas, norm constraints.

\section{Half-spaces and polyhedral sets}
\label{sec:halfspaces-polyhedra}
% Linear inequality constraints. Intersections.

\section{Degenerate sets: EmptySet and FullSpace}
\label{sec:degenerate-sets}

\section{Set operations}
\label{sec:set-operations}
% Intersection, Minkowski sums (via support functions).


% ============================================================================
\chapter{Subspaces}
\label{ch:subspaces}
% ============================================================================

\section{Motivation: equality constraints}
\label{sec:subspace-motivation}
% Linear constraints from conservation laws, boundary conditions.

\section{Orthogonal projectors}
\label{sec:projectors}
% Definition, construction from basis, properties P² = P = P*.

\begin{codebox}
\texttt{OrthogonalProjector.from\_basis(space, [v1, v2, ...])} constructs
the projection $P = \sum_i \hat{v}_i \hat{v}_i^\top$ from spanning vectors.
\end{codebox}

\section{Linear subspaces}
\label{sec:linear-subspaces}

\section{Affine subspaces}
\label{sec:affine-subspaces}
% Translation of linear subspace. Solution set of Bu = w.

\subsection{Construction from tangent bases}
\label{ssec:tangent-basis}

\subsection{Construction from linear equations}
\label{ssec:from-linear-equation}

\section{Tangent spaces and dimension}
\label{sec:tangent-spaces}
% get_tangent_basis(), tolerant Gram--Schmidt.


% ============================================================================
\chapter{Convex Analysis and Support Functions}
\label{ch:convex-analysis}
% ============================================================================

\section{Why convex analysis?}
\label{sec:why-convex-analysis}
% DLI framework: what can we learn without probability?
% Support functions are the key tool.

\begin{whybox}
Convex analysis provides a complete, probability-free framework for
answering the question: given my data and prior constraints, what can I
confidently say about the model? The support function is the central
object that makes this possible.
\end{whybox}

\section{Support functions: intuition}
\label{sec:support-intuition}
% Geometric picture: σ_C(q) = how far C extends in direction q.

\section{Formal definition and properties}
\label{sec:support-definition}
% Convexity, positive homogeneity, subadditivity. Conjugacy.

\section{Subgradients and support points}
\label{sec:subgradients}
% ∂σ_C(q) = argmax, geometric interpretation.

\section{Concrete support functions}
\label{sec:concrete-support}

\subsection{BallSupportFunction}
\label{ssec:ball-support}

\subsection{EllipsoidSupportFunction}
\label{ssec:ellipsoid-support}

\subsection{PolyhedralSupportFunction}
\label{ssec:polyhedral-support}

\subsection{SobolevBallSupportFunction}
\label{ssec:sobolev-support}

\section{Combinators: building complex sets from simple ones}
\label{sec:support-combinators}
% Infimal convolution, sum, scaled sum.

\subsection{Infimal convolution (Minkowski sum)}
\label{ssec:infimal-convolution}

\subsection{Support function sum (intersection)}
\label{ssec:support-sum}

\subsection{Scaled sums}
\label{ssec:scaled-sum}

\section{The dual master equation}
\label{sec:dual-master}
% h_U(q) = inf_λ { ⟨λ, d̃⟩ + σ_B(T'q − G'λ) + σ_V(−λ) }
% Full derivation, geometric interpretation.
% Connection to Al-Attar \& Crawford (2021).


% ============================================================================
\chapter{Gaussian Measures}
\label{ch:gaussian-measures}
% ============================================================================

\section{Why Gaussian measures?}
\label{sec:why-gaussian}
% The natural probabilistic building block for linear-Gaussian problems.

\begin{whybox}
If your prior knowledge about the model is naturally described by a mean
and covariance, and your forward operator is linear, and your measurement
noise is Gaussian, then the full posterior distribution is Gaussian and
can be computed in closed form. This is what \texttt{GaussianMeasure}
provides.
\end{whybox}

\section{Gaussian measures on Hilbert spaces}
\label{sec:gaussian-definition}
% Mean, covariance operator. Cameron--Martin space.
% Difference from finite-dimensional: covariance is an operator.

\section{Sampling}
\label{sec:gaussian-sampling}
% Cholesky-based and randomised approaches.

\section{Push-forward and marginals}
\label{sec:push-forward}
% Am ~ N(Aμ, ACA*).

\section{Conditioning (Bayesian update)}
\label{sec:gaussian-conditioning}
% Posterior formula. Full derivation.


% %%%%%%%%%%%%%%%%%%%%%%%%%%%%%%%%%%%%%%%%%%%%%%%%%%%%%%%%%%%%%%%%%%%%%%%%%%%%
\part{Problems and Algorithms}
\label{part:algorithms}
% %%%%%%%%%%%%%%%%%%%%%%%%%%%%%%%%%%%%%%%%%%%%%%%%%%%%%%%%%%%%%%%%%%%%%%%%%%%%

% ============================================================================
\chapter{Forward Problems}
\label{ch:forward-problems}
% ============================================================================

\section{The forward modelling framework}
\label{sec:forward-framework}
% d = Gm + ε. The forward problem as a data structure.

\section{ForwardProblem and LinearForwardProblem}
\label{sec:forward-classes}
% Constructor, properties, validation.

\begin{codebox}
\texttt{LinearForwardProblem(model\_space, data\_space, forward\_operator,
data\_error\_measure)} bundles the complete forward problem specification.
\end{codebox}


% ============================================================================
\chapter{Deterministic Inversion}
\label{ch:deterministic-inversion}
% ============================================================================

\section{Least-squares inversion}
\label{sec:least-squares}
% min ‖Gm − d‖². Normal equations. Uniqueness.

\section{Minimum-norm inversion}
\label{sec:minimum-norm}
% min ‖m‖ s.t. Gm = d. Pseudo-inverse. Why the norm matters.

\begin{whybox}
The minimum-norm solution depends on the inner product: change the norm,
change the answer. This is why the careful treatment of Hilbert spaces
in \cref{ch:hilbert-spaces} pays off here.
\end{whybox}

\section{Constrained inversions}
\label{sec:constrained-inversions}
% Adding affine constraints Bu = w.

\section{Connection to the dual master equation}
\label{sec:det-inv-dli}
% How deterministic inversions are special cases of DLI.


% ============================================================================
\chapter{Bayesian Inversion}
\label{ch:bayesian-inversion}
% ============================================================================

\section{The Bayesian framework for inverse problems}
\label{sec:bayesian-framework}
% Prior → forward model → likelihood → posterior.
% Stuart (2010) well-posedness.

\section{Linear Bayesian inversion}
\label{sec:linear-bayesian}
% Closed-form posterior. Kalman gain formula.

\begin{codebox}
\texttt{LinearBayesianInversion(forward\_problem, prior)} computes the
posterior Gaussian measure. Access via \texttt{.invert(data)}.
\end{codebox}

\section{Constrained Bayesian inversion}
\label{sec:constrained-bayesian}
% Hard constraints on affine subspaces. Conditional Gaussian.

\section{Posterior uncertainty}
\label{sec:posterior-uncertainty}
% Posterior covariance, marginal variances, confidence ellipsoids.


% ============================================================================
\chapter{Backus--Gilbert and DLI Inference}
\label{ch:backus-gilbert}
% ============================================================================

\section{The inference problem}
\label{sec:inference-problem}
% We want properties p = Tm, not the full model m.

\section{Backus--Gilbert estimators}
\label{sec:bg-estimators}
% Averaging kernels. Resolution--variance trade-off.

\section{The admissible property set $\mathcal{U}$}
\label{sec:admissible-set}
% All property values consistent with data and constraints.
% Characterised by h_U(q).

\section{Computing $\mathcal{U}$ via the dual master equation}
\label{sec:computing-U}
% Numerical evaluation of h_U. Gradient-based optimisation.
% Support points as extreme models.

\section{Polytope outer approximation}
\label{sec:polytope-approximation}
% Three operational regimes: underdetermined, critical, overdetermined.


% ============================================================================
\chapter{Convex Optimisation}
\label{ch:convex-optimisation}
% ============================================================================

\section{Overview of optimisation in pygeoinf}
\label{sec:optim-overview}
% What problems arise, which solvers handle them.

\section{Linear solvers}
\label{sec:linear-solvers}
% CG, MINRES, direct methods.

\section{Convex optimisation algorithms}
\label{sec:convex-algorithms}
% Bundle methods, proximal methods.

\section{Preconditioners}
\label{sec:preconditioners}
% Why preconditioning matters. Available preconditioners.


% %%%%%%%%%%%%%%%%%%%%%%%%%%%%%%%%%%%%%%%%%%%%%%%%%%%%%%%%%%%%%%%%%%%%%%%%%%%%
\part{Computational Methods and Visualisation}
\label{part:computation}
% %%%%%%%%%%%%%%%%%%%%%%%%%%%%%%%%%%%%%%%%%%%%%%%%%%%%%%%%%%%%%%%%%%%%%%%%%%%%

% ============================================================================
\chapter{Computational Utilities}
\label{ch:computational}
% ============================================================================

\section{Randomised matrix decompositions}
\label{sec:randomised}
% Randomised SVD, range finder, Cholesky.

\section{Direct sum operators}
\label{sec:direct-sum-operators}
% Block-diagonal structures.

\section{Parallel computation}
\label{sec:parallel}

\section{Auxiliary functions}
\label{sec:auxiliary}


% ============================================================================
\chapter{Visualisation}
\label{ch:visualisation}
% ============================================================================

\section{The SubspaceSlicePlotter}
\label{sec:slice-plotter}
% Visualising convex sets by slicing.

\section{1D slices}
\label{sec:1d-slices}

\section{2D slices}
\label{sec:2d-slices}

\section{3D slices}
\label{sec:3d-slices}

\section{Worked example: visualising the admissible property set}
\label{sec:vis-example}


% ============================================================================
\appendix
% ============================================================================

% %%%%%%%%%%%%%%%%%%%%%%%%%%%%%%%%%%%%%%%%%%%%%%%%%%%%%%%%%%%%%%%%%%%%%%%%%%%%
\part*{Appendices}
\addcontentsline{toc}{part}{Appendices}
% %%%%%%%%%%%%%%%%%%%%%%%%%%%%%%%%%%%%%%%%%%%%%%%%%%%%%%%%%%%%%%%%%%%%%%%%%%%%

% ============================================================================
\chapter{Convex Analysis Primer}
\label{app:convex-primer}
% ============================================================================

\section{Convex sets and functions}
\label{sec:convex-sets-functions}

\section{Separation theorems}
\label{sec:separation}

\section{Conjugate (Fenchel--Legendre) duality}
\label{sec:fenchel-duality}

\section{Subdifferentials}
\label{sec:subdifferentials}

\section{Infimal convolution}
\label{sec:infimal-conv-appendix}


% ============================================================================
\chapter{Functional Analysis Prerequisites}
\label{app:functional-analysis}
% ============================================================================

\section{Banach spaces}
\label{sec:banach-spaces}

\section{Weak and weak-* topologies}
\label{sec:weak-topologies}

\section{Compact operators}
\label{sec:compact-operators}

\section{Spectral theory for self-adjoint operators}
\label{sec:spectral-theory}


% ============================================================================
\chapter{Deferred Proofs}
\label{app:proofs}
% ============================================================================
% Long proofs that would interrupt the main narrative.


% ============================================================================
\chapter{Complete API Quick Reference}
\label{app:api-reference}
% ============================================================================

\section{hilbert\_space module}
\label{sec:api-hilbert}

\section{linear\_forms module}
\label{sec:api-forms}

\section{linear\_operators module}
\label{sec:api-operators}

\section{subsets module}
\label{sec:api-subsets}

\section{subspaces module}
\label{sec:api-subspaces}

\section{convex\_analysis module}
\label{sec:api-convex}

\section{gaussian\_measure module}
\label{sec:api-gaussian}

\section{forward\_problem module}
\label{sec:api-forward}

\section{inversion module}
\label{sec:api-inversion}

\section{linear\_bayesian module}
\label{sec:api-bayesian}

\section{linear\_optimisation module}
\label{sec:api-optimisation}

\section{backus\_gilbert module}
\label{sec:api-bg}

\section{visualisation module}
\label{sec:api-vis}


% ============================================================================
\backmatter
% ============================================================================

% ---------------------------------------------------------------------------
% Bibliography
% ---------------------------------------------------------------------------
\bibliographystyle{plainnat}
\bibliography{bibliography}

% ---------------------------------------------------------------------------
% Index
% ---------------------------------------------------------------------------
% \printindex  % uncomment when index entries are added

\end{document}
